\documentclass[11pt]{article}
\usepackage[a4paper,margin=25mm]{geometry}
\usepackage{amsmath,amssymb,amsthm,booktabs}
\usepackage[british]{babel}
\usepackage{url,xcolor,graphicx}
\usepackage[round,authoryear]{natbib}
\usepackage[hidelinks]{hyperref}
\hypersetup{pdfauthor={Amanjeet Singh},pdftitle={Capped GBM pricing}}
\setlength{\bibhang}{1.5em}
\setlength{\bibsep}{2pt}
\setlength{\emergencystretch}{1.5em}
\renewcommand{\bibfont}{\small}
\renewcommand{\topfraction}{0.85}
\renewcommand{\textfraction}{0.1}
\renewcommand{\floatpagefraction}{0.8}
\newtheorem{theorem}{Theorem}[section]
\newtheorem{proposition}[theorem]{Proposition}
\title{Capped GBM pricing\\[2mm]\large Replication, numerical convergence and parameter uncertainty}
\author{Amanjeet Singh}
\date{28 October 2026}
\begin{document}
\maketitle
\begin{abstract}
\noindent
The analytic call-spread representation of a capped European claim links
no-arbitrage replication, numerical pricing and parameter uncertainty under
geometric Brownian motion. I derive prices and five Greeks, prove stability
and convergence of a monotone heat-coordinate backward-Euler scheme for the
kinked payoff in exact arithmetic, and compare analytic, PDE and Monte Carlo
calculations. Explicit domain, interpolation and bump-error bounds separate
pricing and sensitivity errors. An $L^2$ projection argument links jointly conditioned controls to variance
reduction, with the variance penalty from independent-pilot coefficients quantified explicitly. Gaussian-return inference distinguishes physical drift
from risk-neutral valuation. Because capped prices can have an interior
volatility maximum, endpoint-only propagation can omit admissible prices.
I derive the maximum and propagate parameter intervals through the complete
price image using exact scatter, bounded observation rounding and outward
evaluation of quantiles and prices. Synthetic experiments retain adverse control
rankings and coarse-quadrature failures. The analytic call spread remains
the strongest known-parameter benchmark.
\end{abstract}
\noindent\begin{minipage}{\linewidth}
\small\emergencystretch=2em
\textit{Keywords:} derivative pricing; geometric Brownian motion; replication;
finite differences; numerical stability; convergence; Monte Carlo Greeks;
conditional expectation; control variates; statistical calibration.\par
\smallskip\textit{JEL classification:} C15, C63, G13.\par
\end{minipage}
% The parent document supplies amsmath, amssymb, amsthm, natbib and hyperref.
% This fragment deliberately contains no claim of theoretical originality.

\section{Financial model and no-arbitrage valuation}
\label{sec:cgbm-model}

I use a one-asset claim to connect replication, martingale valuation,
parabolic equations and statistical estimation under one financial model.
The purpose is to demonstrate the mathematics underlying the numerical
comparisons. The classical valuation framework follows
\citet{blackscholes1973} and \citet{merton1973}; these established methods
do not establish a quantum advantage or a new financial theory. Replication and call valuation specialise
\citet{blackscholes1973}; the continuous-dividend treatment follows
\citet{merton1973}. The pricing and sensitivity derivations below retain
the assumptions required for these specialisations.

Let $S>0$, $K>0$, $H>0$, $\sigma>0$ and $\tau>0$. The interest rate $r$ and
continuous dividend yield $q$ are constant. Write $J=K+H$ and
\begin{equation}
 F(x)=\min\{(x-K)^+,H\}=(x-K)^+-(x-J)^+,
 \qquad D=\exp(-r\tau),\qquad v=\sigma\sqrt{\tau}.
 \label{eq:cgbm-payoff}
\end{equation}
Under the selected pricing measure $Q$, the ex-dividend asset follows
\begin{equation}
 dS_t=(r-q)S_t\,dt+\sigma S_t\,dW_t^Q,
 \qquad
 X=S\exp\{b\tau+vZ\},
 \qquad b=r-q-\tfrac12\sigma^2,
 \label{eq:cgbm-law}
\end{equation}
where $Z$ is standard normal. The price is
\begin{equation}
 V(S,\tau;r,q,\sigma)=D\,\mathbb E_Q[F(X)].
 \label{eq:cgbm-price}
\end{equation}
This is a model comparison with synthetic parameters, rather than a
calibrated empirical test. A physical return mean cannot be substituted
for the pricing drift $r-q$.
I assume frictionless continuous trading, reinvestable dividends,
borrowing and lending at the same rate $r$, and an unrestricted stock
position within the stated admissible strategy class. The probability
space carries the usual augmented Brownian filtration, satisfying the
usual completeness and right-continuity conditions.

\begin{proposition}[Replication and the pricing equation]
\label{prop:cgbm-replication}
The price in \eqref{eq:cgbm-price} is bounded, continuous at expiry and smooth
for $S,\tau>0$. It satisfies
\begin{equation}
 V_\tau=\tfrac12\sigma^2S^2V_{SS}+(r-q)SV_S-rV,
 \qquad V(S,0)=F(S).
 \label{eq:cgbm-pde}
\end{equation}
Holding $V_S$ shares and placing $V-SV_S$ in the money-market account
replicates the claim when dividends are reinvested within the
self-financing strategy. Its discounted wealth is non-negative, and its
stochastic gain is square-integrable on each finite horizon. Under the
usual admissibility condition that discounted wealth is bounded below,
the stipulated market has no arbitrage. This replicating price is the
unique no-arbitrage price among bounded claim valuations and the unique
bounded classical solution of the terminal-value problem. Here a classical
solution is $C^{1,2}$ for positive remaining maturity and is continuous
at expiry.
\end{proposition}

The dividend term belongs in the self-financing gain equation, not in
the terminal payoff. Appendix~\ref{app:cgbm-financial} derives both the
measure change and the replication argument, including the limiting
argument at the payoff kinks.

\section{Analytic prices and Greeks}
\label{sec:cgbm-analytic}

Let $\Phi$ and $\varphi$ denote the standard normal distribution function
and density. For strike $A>0$, define
\begin{equation}
 d_{1,A}=\frac{\log(S/A)+(r-q+\sigma^2/2)\tau}{v},
 \qquad d_{2,A}=d_{1,A}-v,
 \qquad C_A=S e^{-q\tau}\Phi(d_{1,A})-A D\Phi(d_{2,A}).
 \label{eq:cgbm-call}
\end{equation}
Then $V=C_K-C_J$. Write
$\Delta\Phi_i=\Phi(d_{i,K})-\Phi(d_{i,J})$ and
$\Delta\varphi=\varphi(d_{1,K})-\varphi(d_{1,J})$. Differentiation gives
\begin{align}
 \Delta=V_S&=e^{-q\tau}\Delta\Phi_1,\label{eq:cgbm-delta}\\
 \Gamma=V_{SS}&=\frac{e^{-q\tau}\Delta\varphi}{Sv},\label{eq:cgbm-delta-gamma}\\
 \mathcal V=V_\sigma&=S e^{-q\tau}\sqrt{\tau}\,\Delta\varphi,\label{eq:cgbm-vega}\\
 \mathcal R=V_r&=\tau D\{K\Phi(d_{2,K})-J\Phi(d_{2,J})\},\label{eq:cgbm-vega-rho}\\
 \Theta=-V_\tau&=-\frac{S e^{-q\tau}\sigma}{2\sqrt{\tau}}\Delta\varphi
 +q S e^{-q\tau}\Delta\Phi_1\notag\\
 &\quad-rD\{K\Phi(d_{2,K})-J\Phi(d_{2,J})\}.\label{eq:cgbm-theta}
\end{align}
Vega and rho are derivatives per unit change in volatility and interest
rate. Calendar theta uses a fixed expiry and advances the valuation date,
so it is the negative maturity derivative. Gamma and vega can be negative
for a capped claim; this is a property of the call spread, not a numerical
failure.

\section{A monotone heat-equation discretisation}
\label{sec:cgbm-heat}

Introduce $y=\log(S/K)+b\tau$, $a=\sigma^2/2$ and
$f(y)=F(K e^y)$. The whole-line heat solution $w$ satisfies
\begin{equation}
 w_\tau=a w_{yy},\qquad w(0,y)=f(y),\qquad
 V(S,\tau)=D w(\tau,y).
 \label{eq:cgbm-heat}
\end{equation}
Its payoff kinks lie at $0$ and $c=\log(J/K)$. Moreover,
$0\leq f\leq H$ and $f$ is globally Lipschitz with constant
$L_f=J$. For a finite domain $L<0<c<U$, let $w^{\mathrm{box}}$ solve the
same equation with fixed boundary values $0$ and $H$.
These values approximate the exact financial far-field limits
\begin{equation}
 \lim_{S\downarrow0}V(S,\tau)=0,\qquad
 \lim_{S\to\infty}V(S,\tau)=DH.
 \label{eq:cgbm-far-field}
\end{equation}
Indeed the terminal asset tends pointwise to zero or infinity as its
initial spot does, and the bounded continuous payoff permits dominated
convergence. At remaining maturity $\tau$, the heat-domain endpoints
correspond to asset values $K\exp(L-b\tau)$ and
$K\exp(U-b\tau)$. Their finite values do not make the imposed boundary
prices exact; the exit-tail bound below quantifies the approximation.

For integers $M\geq3$ and $N\geq1$, set
$h=(U-L)/M$, $y_j=L+jh$, $k=\tau/N$ and
$\lambda=ak/h^2$. Backward Euler gives
\begin{equation}
 -\lambda W_{j-1}^{n}+(1+2\lambda)W_j^{n}
 -\lambda W_{j+1}^{n}=W_j^{n-1},\qquad
 W_0^n=0,\quad W_M^n=H,\quad W_j^0=f(y_j).
 \label{eq:cgbm-be}
\end{equation}
I interpolate the final grid values linearly at the query $y$ and then
apply the discount factor. Subtracting the affine boundary lift diagonalises
this same finite-difference operator in the discrete sine basis. Its
eigenvalues and backward-Euler factors are
\begin{equation}
 \nu_\ell=\frac{4}{h^2}\sin^2\!\left(\frac{\ell\pi}{2M}\right),
 \qquad (1+ak\nu_\ell)^{-N},\qquad 1\leq\ell\leq M-1.
 \label{eq:cgbm-dst}
\end{equation}
This is a faster evaluation of \eqref{eq:cgbm-be}, rather than a different
spatial method. A direct tridiagonal time-step implementation checks that
equivalence numerically. A floating-point implementation is not an
interval-arithmetic certificate merely because its mathematical scheme
has a proved bound.

\begin{theorem}[Unconditional stability and a nonsmooth-payoff error bound]
\label{thm:cgbm-be}
For every $h,k>0$, \eqref{eq:cgbm-be} preserves $0\leq W_j^n\leq H$.
Differences between two solutions with the same boundary values contract
in the maximum norm at every step. Fix an even, non-negative
$\rho\in C_c^\infty(-1,1)$ with integral one, and set
\begin{equation}
 m_1=\int_{-1}^1|z|\rho(z)\,dz,\qquad
 c_3=\|\rho'''\|_{L^1},\qquad
 0<\eta<\min\{-L,U-c\}.
 \label{eq:cgbm-mollifier}
\end{equation}
For exact arithmetic, the nodal error at maturity obeys
\begin{equation}
 \max_{1\leq j\leq M-1}
 |W_j^N-w^{\mathrm{box}}(\tau,y_j)|
 \leq 2L_fm_1\eta
 +\tau\left(\frac{ah^2}{12}+\frac{a^2k}{2}\right)
 L_f c_3\eta^{-3}.
 \label{eq:cgbm-be-error}
\end{equation}
Consequently the scheme converges to the finite-domain solution, with
a rigorous upper bound of order $(h^2+k)^{1/4}$ for this nonsmooth payoff
on a fixed time interval. This conservative bound does not assert that
the observed rate is exactly one quarter.
\end{theorem}

The proof uses the maximum principle, a compact mollifier and an odd
periodic extension after the boundary lift. It does not assume that the
unmollified payoff has four derivatives. The full proof is in
Appendix~\ref{app:cgbm-numerical}.

\begin{proposition}[Domain and interpolation errors]
\label{prop:cgbm-domain}
For $x\in[L,U]$, put
\begin{equation}
 \beta(\tau,x)=\min\!\left\{1,
 2\Phi\!\left(-\frac{x-L}{\sigma\sqrt{\tau}}\right)
 +2\Phi\!\left(-\frac{U-x}{\sigma\sqrt{\tau}}\right)\right\}.
 \label{eq:cgbm-domain-beta}
\end{equation}
Then $|w(\tau,x)-w^{\mathrm{box}}(\tau,x)|\leq H\beta(\tau,x)$.
Let $y$ lie between grid nodes $y_j,y_{j+1}$ and define
$\beta_* =\max\{\beta(\tau,y_j),\beta(\tau,y_{j+1})\}$.
The interpolated, discounted price $\widehat V$ satisfies
\begin{align}
 |\widehat V-V|\leq D\Bigg[
 &H\beta_*+2L_fm_1\eta
 +\tau\left(\frac{ah^2}{12}+\frac{a^2k}{2}\right)
 L_f c_3\eta^{-3}\notag\\
 &+\frac{h^2 L_f}{8\sigma\sqrt{\tau}}
 \sqrt{\frac{2}{\pi}}\Bigg].
 \label{eq:cgbm-total-error}
\end{align}
\end{proposition}

The four terms control the finite domain, payoff smoothing, numerical
consistency and interpolation. Payoff smoothing is a proof device; the
implemented initial condition remains the original capped payoff.
On a fixed domain, mesh refinement converges to
$w^{\mathrm{box}}$, leaving a domain error relative to the whole-line
financial price. Full-price consistency requires the domain error to
vanish as well.

\section{Greek consistency and Monte Carlo comparisons}
\label{sec:cgbm-greeks}

Nonsmooth payoffs require separate control of sensitivities; the analysis of 
\citet{gilescarter2006} gives a published comparison for Crank--Nicolson/
Rannacher schemes. I prove the bounds for the backward-Euler scheme used here.

I compute PDE Greeks by central price differences on the same domain and
grid. This fixes the numerical heat domain under a bump; the transformed
query $y$, diffusivity $a$ and discount factor still change with the bumped
parameter. The small-bump limit must respect the price error.

\begin{proposition}[Propagation into central differences]
\label{prop:cgbm-bumps}
Suppose $|\widehat V-V|\leq E$ at every required bumped point. For a scalar
parameter $p$, if $|\partial_p^3 V|\leq M_3$ on its bump interval, then
\begin{equation}
 \left|\frac{\widehat V(p+d)-\widehat V(p-d)}{2d}
 -\partial_p V(p)\right|
 \leq \frac{E}{d}+\frac{M_3d^2}{6}.
 \label{eq:cgbm-first-bump}
\end{equation}
For spot gamma, if $|\partial_S^4 V|\leq M_4$ on the spot bump interval,
then
\begin{equation}
 \left|\frac{\widehat V(S+d)-2\widehat V(S)+\widehat V(S-d)}{d^2}
 -\Gamma\right|
 \leq \frac{4E}{d^2}+\frac{M_4d^2}{12}.
 \label{eq:cgbm-second-bump}
\end{equation}
These statements apply uniformly on a compact parameter set with
$S,\sigma,\tau$ bounded away from zero, provided all queries lie inside
the domain. Refining the mesh and enlarging the domain so that $E_n\to0$,
then choosing $d_n\to0$ and $E_n/d_n\to0$ for first derivatives or
$E_n/d_n^2\to0$ for gamma, proves consistency for the whole-line Greeks.
\end{proposition}

After fixing units, choices proportional to $E_n^{1/3}$ and $E_n^{1/4}$
balance the displayed error terms for first derivatives and gamma,
respectively. Practical comparisons also retain the exact analytic
central-difference bias at each implemented bump. This diagnoses numerical
price error separately from bump bias and prevents an apparent Greek
discrepancy from being assigned wholly to the PDE scheme.

The implemented bump recipe uses a dimensionless log-grid spacing $h$,
a positive spot reference scale $S_*$, and positive reference scales for
the other parameters. The spot bump is proportional to $S_*h^{1/8}$;
first-order parameter bumps are proportional to $h^{1/4}$ in their
respective reference units. With $k=O(h^2)$, the rigorous price bound is
$O(h^{1/2})$ after making the domain error $o(h^{1/2})$. The resulting
spot gamma error is $O(h^{1/4})$, and the first-order errors are at most
$O(h^{1/4})$. These are conservative guaranteed upper rates, not fitted
empirical orders. A finite fixed domain retains its displayed exit-tail
allowance, regardless of this bump recipe.

Pathwise and likelihood-ratio estimators have established foundations
\citep{broadieglasserman1996}. The capped-payoff formulae and their ranges
are derived explicitly in Appendix~\ref{app:cgbm-monte-carlo}.

For exact terminal sampling, write $I=\mathbf 1_{\{K<X<J\}}$ and
define the following one-path estimators:
\begin{align}
 Y_V&=D F(X),
 &Y_\Delta&=D I X/S,\notag\\
 Y_{\mathcal V}&=D I X(\sqrt{\tau}Z-\sigma\tau),
 &Y_{\mathcal R}&=D\tau\{IX-F(X)\},
 \label{eq:cgbm-pw-estimators}\\
 Y_\Theta&=D\left[rF(X)-IX
 \left(b+\frac{\sigma Z}{2\sqrt{\tau}}\right)\right],
 \label{eq:cgbm-theta-estimator}\\
 Y_\Gamma&=\frac{D F(X)}{S^2}
 \left\{\frac{Z^2-1}{v^2}-\frac{Z}{v}\right\}.
 \label{eq:cgbm-lr-estimator}
\end{align}
The first-order Greeks use pathwise differentiation, following the
simulation sensitivity framework of \citet{broadieglasserman1996}.
Gamma uses the
second likelihood-ratio weight: differentiating the pathwise indicator
as though it were smooth would omit its boundary contribution.
I also retain the bounded mixed pathwise/likelihood-ratio rival
\begin{equation}
 Y_{\Gamma,\mathrm{mix}}
 =\frac{D I X}{S^2}\left(\frac{Z}{v}-1\right).
 \label{eq:cgbm-mixed-gamma}
\end{equation}
It follows by differentiating delta's fixed-asset density representation,
rather than differentiating a random indicator pointwise. Its derivation
appears in Appendix~\ref{app:cgbm-monte-carlo}. Boundedness gives a
different finite-sample uncertainty guarantee, but does not by itself
prove variance dominance over the ordinary likelihood-ratio estimator.

\begin{proposition}[Unbiased estimators and finite-sample uncertainty]
\label{prop:cgbm-mc}
Each estimator in \eqref{eq:cgbm-pw-estimators}--\eqref{eq:cgbm-mixed-gamma}
has expectation equal to its analytic price or Greek and finite second
moment. Let
\begin{equation}
 z_- =\frac{\log(K/S)-b\tau}{v},\quad
 z_+ =\frac{\log(J/S)-b\tau}{v},\quad
 m=\max\{|z_-|,|z_+|\}.
 \label{eq:cgbm-active-z}
\end{equation}
Deterministic ranges for price, delta, vega, rho and theta are, respectively,
\begin{align}
 &[0,DH],\quad[0,DJ/S],\quad[-B_{\mathcal V},B_{\mathcal V}],
 \quad[-D\tau H,D\tau K],\quad[-B_\Theta,B_\Theta],\notag\\
 &B_{\mathcal V}=DJ(\sqrt{\tau}m+\sigma\tau),\qquad
 B_\Theta=D\left[|r|H+J\left(|b|+
 \frac{\sigma m}{2\sqrt{\tau}}\right)\right].
 \label{eq:cgbm-mc-ranges}
\end{align}
For $n$ independent raw paths and a range width $R$, a distribution-free
two-sided confidence radius with failure probability at most $\alpha$,
following \citet{hoeffding1963}, is
\begin{equation}
 R\sqrt{\frac{\log(2/\alpha)}{2n}}.
 \label{eq:cgbm-hoeffding}
\end{equation}
For gamma, the bound
\begin{equation}
 \operatorname{Var}(Y_\Gamma)
 \leq B_\Gamma^2
 :=\frac{D^2H^2}{S^4}\left(\frac{2}{v^4}+\frac{1}{v^2}\right)
 \label{eq:cgbm-gamma-variance}
\end{equation}
gives the finite-sample Chebyshev radius $B_\Gamma/\sqrt{n\alpha}$.
The mixed gamma estimator has range
\begin{equation}
 [-B_{\mathrm{mix}},B_{\mathrm{mix}}],\qquad
 B_{\mathrm{mix}}=\frac{DJ}{S^2}\left(\frac{m}{v}+1\right),
 \label{eq:cgbm-mixed-gamma-bound}
\end{equation}
so it admits the Hoeffding radius with width $2B_{\mathrm{mix}}$.
Allocating failure probabilities across outputs and using a union bound
gives a simultaneous statement. Sample-variance normal intervals are
separate asymptotic diagnostics and do not have these finite-sample
coverage guarantees.
\end{proposition}

The numerical study allocates $0.05/49$ to each of the six ordinary raw
outputs and the additional raw mixed gamma across seven contracts.
The corresponding 49-interval family has failure probability at most
0.05 under the ideal independent-normal sampling model. Comparisons
between Monte Carlo methods additionally report sample-variance normal
intervals as asymptotic diagnostics.

The strongest implemented Monte Carlo comparison combines antithetic
normal pairs with a centred discounted-stock control,
$D X-S e^{-q\tau}$. I fit each control coefficient on independent pilot
pairs and hold it fixed for the main sample. Conditional on the pilot,
the controlled mean is unbiased because the control has expectation zero.
Each independent pair, rather than its two correlated paths, is the
sampling unit. The pilot and both terminal evaluations per pair belong
in the computational work. The controlled estimators contain an unbounded
lognormal control, so their sample-variance intervals remain asymptotic;
the bounded raw-path Hoeffding guarantees are not transferred to them.

\section{Synthetic comparison design and numerical evidence}
\label{sec:cgbm-results}

I use seven designed contracts, all with $K=100$, to exercise both payoff
kinks, the interior spread, deep capping, a short low-volatility horizon,
negative interest with dividends and a volatile long horizon.
Table~\ref{tab:cgbm-contracts} defines the inputs. No observations are
estimated from LSEG or another data provider. The full comparisons and
source/code hashes are retained in \texttt{results/capped\_gbm}.

\begin{table}[tb]\centering\small
\begin{tabular}{lrrrrrr}\toprule
Contract & $S$ & $H$ & $\tau$ & $\sigma$ & $r$ & $q$\\\midrule
Lower kink & 100 & 20 & 1 & 0.2 & 0.03 & 0.01 \\
Upper kink & 120 & 20 & 1 & 0.2 & 0.03 & 0.01 \\
Inside spread & 110 & 20 & 1 & 0.2 & 0.03 & 0.01 \\
Deep cap & 160 & 20 & 1 & 0.2 & 0.03 & 0.01 \\
Low-volatility short dated & 100 & 20 & 0.05 & 0.06 & 0.03 & 0.01 \\
Negative rate/dividend & 105 & 20 & 1 & 0.25 & -0.015 & 0.045 \\
High volatility & 95 & 30 & 2 & 0.65 & 0.03 & 0.01 \\
\bottomrule\end{tabular}
\caption{Wholly synthetic inputs. Rates, volatility and dividend yields are expressed per year; the cap and spot use common monetary units. The labels describe designed test cases rather than an empirical asset sample.}
\label{tab:cgbm-contracts}\end{table}

The coupled meshes use $M=128,256,512,1024,2048,4096,8192,16384$ and
$N=\lceil M^2/8\rceil$. Additional experiments refine time at fixed space
and space at fixed time. The per-contract heat boxes contain both kinks
and use nine diffusion standard deviations of margin, with a minimum
margin of $0.3$. These finite experiments keep the box fixed; the
whole-line convergence theorem additionally requires its tail allowance
to vanish. Domain-enlargement checks preserve both $h$ and $k$, so a
boundary change is not confused with changing the spatial resolution.

The spot bump is
$\min\{0.1S,S(0.006h^{1/8}+h)\}$.
Volatility, rate and maturity bumps are respectively
$\min\{0.25\sigma,0.10\sigma h^{1/4}\}$,
$0.01h^{1/4}$ and
$\min\{0.25\tau,0.05\tau h^{1/4}\}$.
All actual steps and signed bump biases are retained per comparison.
The analytic central difference uses precisely the PDE bump and therefore
isolates its bias. The tridiagonal and sine-basis solvers are compared
at one, three and 64 time steps, including large values of $\lambda$.

\begin{table}[tb]\centering\small
\begin{tabular}{lrl}\toprule
Quantity & Maximum absolute error at $M=16384$ & Units\\\midrule
Price & $9.950\!\times\!10^{-7}$ & Monetary units \\
Delta & $1.111\!\times\!10^{-4}$ & Per unit spot \\
Gamma & $4.137\!\times\!10^{-4}$ & Per unit spot squared \\
Vega & $1.829\!\times\!10^{-3}$ & Per unit volatility \\
Rho & $3.171\!\times\!10^{-4}$ & Per unit rate \\
Theta & $1.026\!\times\!10^{-3}$ & Per year of calendar time \\
\bottomrule\end{tabular}
\caption{Worst absolute discrepancy from the analytic price/Greek across the seven synthetic contracts. Greek errors include the declared bump bias. Different derivative units are not combined into a single accuracy score.}
\label{tab:cgbm-fine-errors}\end{table}

\begin{figure}[tb]
\centering
\setlength{\unitlength}{1mm}
\resizebox{\linewidth}{!}{%
\begin{picture}(164,71)%
\put(14,62){\makebox(61,5)[c]{\small (a) Absolute price error}}%
\put(14.000,22){\line(0,1){38}}%
\put(14.000,17){\makebox(0,3){\scriptsize 128}}%
{\color[HTML]{E5E5E5}\put(22.714,22){\line(0,1){38}}}%
{\color[HTML]{E5E5E5}\put(31.429,22){\line(0,1){38}}}%
\put(31.429,17){\makebox(0,3){\scriptsize 512}}%
{\color[HTML]{E5E5E5}\put(40.143,22){\line(0,1){38}}}%
{\color[HTML]{E5E5E5}\put(48.857,22){\line(0,1){38}}}%
\put(48.857,17){\makebox(0,3){\scriptsize 2048}}%
{\color[HTML]{E5E5E5}\put(57.571,22){\line(0,1){38}}}%
{\color[HTML]{E5E5E5}\put(66.286,22){\line(0,1){38}}}%
\put(66.286,17){\makebox(0,3){\scriptsize 8192}}%
\put(75.000,22){\line(0,1){38}}%
\put(75.000,17){\makebox(0,3){\scriptsize 16384}}%
{\color[HTML]{E5E5E5}\put(14.000,22.000){\line(1,0){61}}}%
\put(12.000,20.500){\makebox(0,3)[r]{\scriptsize $10^{-8}$}}%
{\color[HTML]{E5E5E5}\put(14.000,28.333){\line(1,0){61}}}%
\put(12.000,26.833){\makebox(0,3)[r]{\scriptsize $10^{-7}$}}%
{\color[HTML]{E5E5E5}\put(14.000,34.667){\line(1,0){61}}}%
\put(12.000,33.167){\makebox(0,3)[r]{\scriptsize $10^{-6}$}}%
{\color[HTML]{E5E5E5}\put(14.000,41.000){\line(1,0){61}}}%
\put(12.000,39.500){\makebox(0,3)[r]{\scriptsize $10^{-5}$}}%
{\color[HTML]{E5E5E5}\put(14.000,47.333){\line(1,0){61}}}%
\put(12.000,45.833){\makebox(0,3)[r]{\scriptsize $10^{-4}$}}%
{\color[HTML]{E5E5E5}\put(14.000,53.667){\line(1,0){61}}}%
\put(12.000,52.167){\makebox(0,3)[r]{\scriptsize $10^{-3}$}}%
{\color[HTML]{E5E5E5}\put(14.000,60.000){\line(1,0){61}}}%
\put(12.000,58.500){\makebox(0,3)[r]{\scriptsize $10^{-2}$}}%
\put(14,13){\makebox(61,3)[c]{\scriptsize Space intervals $M$}}%
{\color[HTML]{0072B2}\linethickness{0.5pt}\qbezier(14.000,58.723)(18.357,55.865)(22.714,53.006)}%
{\color[HTML]{0072B2}\linethickness{0.5pt}\qbezier(22.714,53.006)(27.071,49.345)(31.429,45.683)}%
{\color[HTML]{0072B2}\linethickness{0.5pt}\qbezier(31.429,45.683)(35.786,42.675)(40.143,39.667)}%
{\color[HTML]{0072B2}\linethickness{0.5pt}\qbezier(40.143,39.667)(44.500,39.188)(48.857,38.708)}%
{\color[HTML]{0072B2}\linethickness{0.5pt}\qbezier(48.857,38.708)(53.214,35.349)(57.571,31.990)}%
{\color[HTML]{0072B2}\linethickness{0.5pt}\qbezier(57.571,31.990)(61.929,32.270)(66.286,32.551)}%
{\color[HTML]{0072B2}\linethickness{0.5pt}\qbezier(66.286,32.551)(70.643,32.055)(75.000,31.560)}%
{\color[HTML]{0072B2}\put(14.000,58.723){\circle*{0.8}}}%
{\color[HTML]{0072B2}\put(22.714,53.006){\circle*{0.8}}}%
{\color[HTML]{0072B2}\put(31.429,45.683){\circle*{0.8}}}%
{\color[HTML]{0072B2}\put(40.143,39.667){\circle*{0.8}}}%
{\color[HTML]{0072B2}\put(48.857,38.708){\circle*{0.8}}}%
{\color[HTML]{0072B2}\put(57.571,31.990){\circle*{0.8}}}%
{\color[HTML]{0072B2}\put(66.286,32.551){\circle*{0.8}}}%
{\color[HTML]{0072B2}\put(75.000,31.560){\circle*{0.8}}}%
{\color[HTML]{0072B2}\put(14.000,8){\line(1,0){4}}}\put(19.000,6.5){\makebox(0,3)[l]{\scriptsize Lower kink}}%
{\color[HTML]{D55E00}\linethickness{0.5pt}\qbezier(14.000,59.561)(18.357,56.515)(22.714,53.469)}%
{\color[HTML]{D55E00}\linethickness{0.5pt}\qbezier(22.714,53.469)(27.071,51.001)(31.429,48.533)}%
{\color[HTML]{D55E00}\linethickness{0.5pt}\qbezier(31.429,48.533)(35.786,46.170)(40.143,43.808)}%
{\color[HTML]{D55E00}\linethickness{0.5pt}\qbezier(40.143,43.808)(44.500,42.554)(48.857,41.300)}%
{\color[HTML]{D55E00}\linethickness{0.5pt}\qbezier(48.857,41.300)(53.214,37.178)(57.571,33.057)}%
{\color[HTML]{D55E00}\linethickness{0.5pt}\qbezier(57.571,33.057)(61.929,33.850)(66.286,34.644)}%
{\color[HTML]{D55E00}\linethickness{0.5pt}\qbezier(66.286,34.644)(70.643,33.475)(75.000,32.306)}%
{\color[HTML]{D55E00}\put(14.000,59.561){\circle*{0.8}}}%
{\color[HTML]{D55E00}\put(22.714,53.469){\circle*{0.8}}}%
{\color[HTML]{D55E00}\put(31.429,48.533){\circle*{0.8}}}%
{\color[HTML]{D55E00}\put(40.143,43.808){\circle*{0.8}}}%
{\color[HTML]{D55E00}\put(48.857,41.300){\circle*{0.8}}}%
{\color[HTML]{D55E00}\put(57.571,33.057){\circle*{0.8}}}%
{\color[HTML]{D55E00}\put(66.286,34.644){\circle*{0.8}}}%
{\color[HTML]{D55E00}\put(75.000,32.306){\circle*{0.8}}}%
{\color[HTML]{D55E00}\put(46.000,8){\line(1,0){4}}}\put(51.000,6.5){\makebox(0,3)[l]{\scriptsize Upper kink}}%
{\color[HTML]{009E73}\linethickness{0.5pt}\qbezier(14.000,50.759)(18.357,53.165)(22.714,55.572)}%
{\color[HTML]{009E73}\linethickness{0.5pt}\qbezier(22.714,55.572)(27.071,51.773)(31.429,47.974)}%
{\color[HTML]{009E73}\linethickness{0.5pt}\qbezier(31.429,47.974)(35.786,47.221)(40.143,46.469)}%
{\color[HTML]{009E73}\linethickness{0.5pt}\qbezier(40.143,46.469)(44.500,45.572)(48.857,44.675)}%
{\color[HTML]{009E73}\linethickness{0.5pt}\qbezier(48.857,44.675)(53.214,43.009)(57.571,41.342)}%
{\color[HTML]{009E73}\linethickness{0.5pt}\qbezier(57.571,41.342)(61.929,37.264)(66.286,33.186)}%
{\color[HTML]{009E73}\linethickness{0.5pt}\qbezier(66.286,33.186)(70.643,30.244)(75.000,27.302)}%
{\color[HTML]{009E73}\put(14.000,50.759){\circle*{0.8}}}%
{\color[HTML]{009E73}\put(22.714,55.572){\circle*{0.8}}}%
{\color[HTML]{009E73}\put(31.429,47.974){\circle*{0.8}}}%
{\color[HTML]{009E73}\put(40.143,46.469){\circle*{0.8}}}%
{\color[HTML]{009E73}\put(48.857,44.675){\circle*{0.8}}}%
{\color[HTML]{009E73}\put(57.571,41.342){\circle*{0.8}}}%
{\color[HTML]{009E73}\put(66.286,33.186){\circle*{0.8}}}%
{\color[HTML]{009E73}\put(75.000,27.302){\circle*{0.8}}}%
{\color[HTML]{009E73}\put(14.000,4){\line(1,0){4}}}\put(19.000,2.5){\makebox(0,3)[l]{\scriptsize Short dated}}%
{\color[HTML]{CC79A7}\linethickness{0.5pt}\qbezier(14.000,58.452)(18.357,57.875)(22.714,57.298)}%
{\color[HTML]{CC79A7}\linethickness{0.5pt}\qbezier(22.714,57.298)(27.071,54.178)(31.429,51.059)}%
{\color[HTML]{CC79A7}\linethickness{0.5pt}\qbezier(31.429,51.059)(35.786,51.974)(40.143,52.888)}%
{\color[HTML]{CC79A7}\linethickness{0.5pt}\qbezier(40.143,52.888)(44.500,50.223)(48.857,47.557)}%
{\color[HTML]{CC79A7}\linethickness{0.5pt}\qbezier(48.857,47.557)(53.214,42.496)(57.571,37.434)}%
{\color[HTML]{CC79A7}\linethickness{0.5pt}\qbezier(57.571,37.434)(61.929,37.659)(66.286,37.885)}%
{\color[HTML]{CC79A7}\linethickness{0.5pt}\qbezier(66.286,37.885)(70.643,36.269)(75.000,34.653)}%
{\color[HTML]{CC79A7}\put(14.000,58.452){\circle*{0.8}}}%
{\color[HTML]{CC79A7}\put(22.714,57.298){\circle*{0.8}}}%
{\color[HTML]{CC79A7}\put(31.429,51.059){\circle*{0.8}}}%
{\color[HTML]{CC79A7}\put(40.143,52.888){\circle*{0.8}}}%
{\color[HTML]{CC79A7}\put(48.857,47.557){\circle*{0.8}}}%
{\color[HTML]{CC79A7}\put(57.571,37.434){\circle*{0.8}}}%
{\color[HTML]{CC79A7}\put(66.286,37.885){\circle*{0.8}}}%
{\color[HTML]{CC79A7}\put(75.000,34.653){\circle*{0.8}}}%
{\color[HTML]{CC79A7}\put(46.000,4){\line(1,0){4}}}\put(51.000,2.5){\makebox(0,3)[l]{\scriptsize High volatility}}%
\put(100,62){\makebox(61,5)[c]{\small (b) Absolute gamma error}}%
\put(100.000,22){\line(0,1){38}}%
\put(100.000,17){\makebox(0,3){\scriptsize 128}}%
{\color[HTML]{E5E5E5}\put(108.714,22){\line(0,1){38}}}%
{\color[HTML]{E5E5E5}\put(117.429,22){\line(0,1){38}}}%
\put(117.429,17){\makebox(0,3){\scriptsize 512}}%
{\color[HTML]{E5E5E5}\put(126.143,22){\line(0,1){38}}}%
{\color[HTML]{E5E5E5}\put(134.857,22){\line(0,1){38}}}%
\put(134.857,17){\makebox(0,3){\scriptsize 2048}}%
{\color[HTML]{E5E5E5}\put(143.571,22){\line(0,1){38}}}%
{\color[HTML]{E5E5E5}\put(152.286,22){\line(0,1){38}}}%
\put(152.286,17){\makebox(0,3){\scriptsize 8192}}%
\put(161.000,22){\line(0,1){38}}%
\put(161.000,17){\makebox(0,3){\scriptsize 16384}}%
{\color[HTML]{E5E5E5}\put(100.000,22.000){\line(1,0){61}}}%
\put(98.000,20.500){\makebox(0,3)[r]{\scriptsize $10^{-5}$}}%
{\color[HTML]{E5E5E5}\put(100.000,31.500){\line(1,0){61}}}%
\put(98.000,30.000){\makebox(0,3)[r]{\scriptsize $10^{-4}$}}%
{\color[HTML]{E5E5E5}\put(100.000,41.000){\line(1,0){61}}}%
\put(98.000,39.500){\makebox(0,3)[r]{\scriptsize $10^{-3}$}}%
{\color[HTML]{E5E5E5}\put(100.000,50.500){\line(1,0){61}}}%
\put(98.000,49.000){\makebox(0,3)[r]{\scriptsize $10^{-2}$}}%
{\color[HTML]{E5E5E5}\put(100.000,60.000){\line(1,0){61}}}%
\put(98.000,58.500){\makebox(0,3)[r]{\scriptsize $10^{-1}$}}%
\put(100,13){\makebox(61,3)[c]{\scriptsize Space intervals $M$}}%
{\color[HTML]{0072B2}\linethickness{0.5pt}\qbezier(100.000,47.306)(104.357,47.409)(108.714,47.513)}%
{\color[HTML]{0072B2}\linethickness{0.5pt}\qbezier(108.714,47.513)(113.071,46.206)(117.429,44.900)}%
{\color[HTML]{0072B2}\linethickness{0.5pt}\qbezier(117.429,44.900)(121.786,42.675)(126.143,40.449)}%
{\color[HTML]{0072B2}\linethickness{0.5pt}\qbezier(126.143,40.449)(130.500,40.749)(134.857,41.049)}%
{\color[HTML]{0072B2}\linethickness{0.5pt}\qbezier(134.857,41.049)(139.214,40.563)(143.571,40.077)}%
{\color[HTML]{0072B2}\linethickness{0.5pt}\qbezier(143.571,40.077)(147.929,38.847)(152.286,37.617)}%
{\color[HTML]{0072B2}\linethickness{0.5pt}\qbezier(152.286,37.617)(156.643,37.488)(161.000,37.359)}%
{\color[HTML]{0072B2}\put(100.000,47.306){\circle*{0.8}}}%
{\color[HTML]{0072B2}\put(108.714,47.513){\circle*{0.8}}}%
{\color[HTML]{0072B2}\put(117.429,44.900){\circle*{0.8}}}%
{\color[HTML]{0072B2}\put(126.143,40.449){\circle*{0.8}}}%
{\color[HTML]{0072B2}\put(134.857,41.049){\circle*{0.8}}}%
{\color[HTML]{0072B2}\put(143.571,40.077){\circle*{0.8}}}%
{\color[HTML]{0072B2}\put(152.286,37.617){\circle*{0.8}}}%
{\color[HTML]{0072B2}\put(161.000,37.359){\circle*{0.8}}}%
{\color[HTML]{0072B2}\put(100.000,8){\line(1,0){4}}}\put(105.000,6.5){\makebox(0,3)[l]{\scriptsize Total error}}%
{\color[HTML]{D55E00}\linethickness{0.5pt}\qbezier(100.000,50.957)(104.357,49.188)(108.714,47.418)}%
{\color[HTML]{D55E00}\linethickness{0.5pt}\qbezier(108.714,47.418)(113.071,45.986)(117.429,44.555)}%
{\color[HTML]{D55E00}\linethickness{0.5pt}\qbezier(117.429,44.555)(121.786,43.458)(126.143,42.362)}%
{\color[HTML]{D55E00}\linethickness{0.5pt}\qbezier(126.143,42.362)(130.500,41.536)(134.857,40.710)}%
{\color[HTML]{D55E00}\linethickness{0.5pt}\qbezier(134.857,40.710)(139.214,40.072)(143.571,39.434)}%
{\color[HTML]{D55E00}\linethickness{0.5pt}\qbezier(143.571,39.434)(147.929,38.915)(152.286,38.397)}%
{\color[HTML]{D55E00}\linethickness{0.5pt}\qbezier(152.286,38.397)(156.643,37.949)(161.000,37.501)}%
{\color[HTML]{D55E00}\put(100.000,50.957){\circle*{0.8}}}%
{\color[HTML]{D55E00}\put(108.714,47.418){\circle*{0.8}}}%
{\color[HTML]{D55E00}\put(117.429,44.555){\circle*{0.8}}}%
{\color[HTML]{D55E00}\put(126.143,42.362){\circle*{0.8}}}%
{\color[HTML]{D55E00}\put(134.857,40.710){\circle*{0.8}}}%
{\color[HTML]{D55E00}\put(143.571,39.434){\circle*{0.8}}}%
{\color[HTML]{D55E00}\put(152.286,38.397){\circle*{0.8}}}%
{\color[HTML]{D55E00}\put(161.000,37.501){\circle*{0.8}}}%
{\color[HTML]{D55E00}\put(132.000,8){\line(1,0){4}}}\put(137.000,6.5){\makebox(0,3)[l]{\scriptsize Bump bias}}%
{\color[HTML]{009E73}\linethickness{0.5pt}\qbezier(100.000,48.761)(104.357,40.326)(108.714,31.892)}%
{\color[HTML]{009E73}\linethickness{0.5pt}\qbezier(108.714,31.892)(113.071,39.741)(117.429,47.591)}%
{\color[HTML]{009E73}\linethickness{0.5pt}\qbezier(117.429,47.591)(121.786,42.931)(126.143,38.272)}%
{\color[HTML]{009E73}\linethickness{0.5pt}\qbezier(126.143,38.272)(130.500,34.421)(134.857,30.569)}%
{\color[HTML]{009E73}\linethickness{0.5pt}\qbezier(134.857,30.569)(139.214,31.330)(143.571,32.092)}%
{\color[HTML]{009E73}\linethickness{0.5pt}\qbezier(143.571,32.092)(147.929,31.614)(152.286,31.137)}%
{\color[HTML]{009E73}\linethickness{0.5pt}\qbezier(152.286,31.137)(156.643,27.335)(161.000,23.533)}%
{\color[HTML]{009E73}\put(100.000,48.761){\circle*{0.8}}}%
{\color[HTML]{009E73}\put(108.714,31.892){\circle*{0.8}}}%
{\color[HTML]{009E73}\put(117.429,47.591){\circle*{0.8}}}%
{\color[HTML]{009E73}\put(126.143,38.272){\circle*{0.8}}}%
{\color[HTML]{009E73}\put(134.857,30.569){\circle*{0.8}}}%
{\color[HTML]{009E73}\put(143.571,32.092){\circle*{0.8}}}%
{\color[HTML]{009E73}\put(152.286,31.137){\circle*{0.8}}}%
{\color[HTML]{009E73}\put(161.000,23.533){\circle*{0.8}}}%
{\color[HTML]{009E73}\put(100.000,4){\line(1,0){4}}}\put(105.000,2.5){\makebox(0,3)[l]{\scriptsize PDE difference}}%
\end{picture}}%
\caption{Synthetic convergence evidence. Panel (a) contrasts the lower and upper payoff kinks, a short-dated low-volatility claim and a high-volatility claim. Panel (b) separates the short-dated gamma total error into analytic bump bias and the PDE difference at the same bump. Absolute component errors need not add because their signed contributions can cancel. Both axes are logarithmic. All seven contracts and all five Greeks are retained in the comparison files.}
\label{fig:cgbm-convergence}
\end{figure}
Figure~\ref{fig:cgbm-convergence} exposes non-monotone refinement rather
than imposing a fitted rate. Kink/grid alignment and interpolation can
affect price error; differencing amplifies these errors in the Greeks.
For the short-dated gamma at the finest mesh, the analytic bump bias is
approximately $-4.28\times10^{-4}$, while the PDE difference at that bump
is approximately $1.45\times10^{-5}$. The larger gamma discrepancy is
therefore mainly a finite-bump effect. Finer price accuracy alone does not
remove it. Independent mesh refinements and all five Greek comparisons
remain available even where a selected plot would conceal adverse behaviour.

Every Monte Carlo method receives 262,144 terminal-payoff evaluations per
contract. Raw sampling uses that many independent normals. Antithetic
sampling uses 131,072 independent pairs. Controlled sampling uses 16,384
independent pilot pairs and 114,688 production pairs, charging both
terminal evaluations of each pair. A control coefficient is fitted
separately for each output and clipped to $[-10^6,10^6]$; the pilot is
independent of production. Shared production normals between methods are
declared, and no confidence claim is made for their paired differences.
The seed is 732451, with deterministic case offsets.
Terminal-evaluation budgets measure sampling work; arithmetic costs and
wall-clock time are not asserted to be identical between estimators.

The 49 raw outputs comprise the price, five Greeks and additional mixed
gamma in each of seven contracts. Each receives failure allowance
$0.05/49$. Their simultaneous finite-sample bounds apply to ideal
independent exact normal draws and exact arithmetic. Their floating-point
constants and pseudo-random normal generation have not been formally
certified. Normal sample-variance intervals are reported separately as
pointwise asymptotic diagnostics. Observed inclusion in the conservative
bounds is a sanity check, not a proof of coverage.

\begin{table}[tb]\centering\small
\begin{tabular}{lrrrr}\toprule
Contract & Analytic $\Gamma$ & Raw LR SE & Raw mixed SE & Variance ratio\\\midrule
Lower kink & $4.026\!\times\!10^{-3}$ & $7.689\!\times\!10^{-5}$ & $1.966\!\times\!10^{-5}$ & 15.293 \\
Upper kink & $-7.259\!\times\!10^{-3}$ & $6.210\!\times\!10^{-5}$ & $2.269\!\times\!10^{-5}$ & 7.493 \\
Inside spread & $-3.183\!\times\!10^{-3}$ & $7.019\!\times\!10^{-5}$ & $1.523\!\times\!10^{-5}$ & 21.240 \\
Deep cap & $-2.747\!\times\!10^{-3}$ & $4.014\!\times\!10^{-5}$ & $2.021\!\times\!10^{-5}$ & 3.945 \\
Low-volatility short dated & $0.296226$ & $3.320\!\times\!10^{-3}$ & $8.632\!\times\!10^{-4}$ & 14.792 \\
Negative rate/dividend & $2.713\!\times\!10^{-3}$ & $4.478\!\times\!10^{-5}$ & $1.272\!\times\!10^{-5}$ & 12.400 \\
High volatility & $-3.680\!\times\!10^{-4}$ & $4.251\!\times\!10^{-6}$ & $2.287\!\times\!10^{-6}$ & 3.455 \\
\bottomrule\end{tabular}
\caption{Raw gamma estimators at equal terminal-evaluation budgets. SE is the measured standard error, not a finite-sample confidence radius. The variance ratio is ordinary payoff-LR estimator variance divided by mixed-estimator variance. All antithetic and independently fitted control variants remain in the full comparison.}
\label{tab:cgbm-gamma-mc}\end{table}

The bounded mixed estimator reduces measured raw gamma variance in every
tested contract. This is an observed comparison, not uniform dominance
over model parameters. Controlled ordinary likelihood-ratio gamma remains
more efficient in several cases. Stock controls also worsen measured
delta and rho efficiency for the interior-spread contract at the matched
total budget. The analytic call spread has neither sampling error nor a
finite-domain or mesh error and remains the strongest baseline here.

\section{Conditional valuation and jointly transformed controls}
\label{sec:cgbm-projection}

I use conditional valuation to connect measure-theoretic probability with
an implemented variance comparison. Conditional Monte Carlo is established
\citep{asmussen2018}. The following argument supplies complete proofs and
an explicit GBM specialisation; it does not claim a new quantum theory.
The expectation and variance in this section are under the stipulated
pricing measure $Q$. For a sub-$\sigma$-algebra $\mathcal G$, write
$P_{\mathcal G}Y=\mathbb E_Q[Y\mid\mathcal G]$. The space
$L^2(\mathcal G)$ consists of equivalence classes admitting a
$\mathcal G$-measurable square-integrable representative.

\begin{theorem}[Conditional expectation as an orthogonal projection]
\label{thm:cgbm-projection}
For $Y\in L^2(Q)$, $L^2(\mathcal G)$ is a closed subspace of $L^2(Q)$
and $m=P_{\mathcal G}Y$ is its unique nearest element. Equivalently,
$m\in L^2(\mathcal G)$ and
$\mathbb E_Q[\mathbf1_A m]=\mathbb E_Q[\mathbf1_A Y]$
for every $A\in\mathcal G$. For every $Z\in L^2(\mathcal G)$,
\begin{align}
 \mathbb E_Q[(Y-Z)^2]
 &=\mathbb E_Q[(Y-m)^2]+\mathbb E_Q[(m-Z)^2],
 \label{eq:cgbm-projection-pythagoras}\\
 m^2&\leq\mathbb E_Q[Y^2\mid\mathcal G]\quad Q\text{-almost surely},
 \label{eq:cgbm-conditional-jensen}\\
 \operatorname{Var}_Q(Y)
 &=\operatorname{Var}_Q(m)+\mathbb E_Q[(Y-m)^2].
 \label{eq:cgbm-projection-total-variance}
\end{align}
For $\mathcal G_1\subseteq\mathcal G_2$, the projections satisfy the
tower property. The variance of the conditional mean increases with
information, while its mean-square prediction error decreases.
\end{theorem}

The two variance statements concern different objects. Replacing $Y$ by
$m$ in an estimator lowers its intrinsic sampling variance. Increasing
the retained information makes $m$ closer to $Y$ and consequently raises
the variance of $m$. Appendix~\ref{app:cgbm-projection} proves the closedness,
projection characterisation, conditional inequality and nested-information
identities without assuming the result being demonstrated.

The control-variate construction is established \citep{nelson1990}.
For a centred control vector $C\in L^2(Q)^d$, let
$\Sigma_C=\mathbb E_Q[CC^{\mathsf T}]$,
$\kappa_C=\mathbb E_Q[C(Y-\mathbb E_QY)]$ and
$\mathcal J(\beta)=\operatorname{Var}_Q(Y-\beta^{\mathsf T}C)$.
Let $C_{\mathcal G}=P_{\mathcal G}C$, componentwise, and define
$\mathcal J_{\mathcal G}(\beta)=
\operatorname{Var}_Q(m-\beta^{\mathsf T}C_{\mathcal G})$.

\begin{proposition}[Conditioning and control optimisation]
\label{prop:cgbm-projected-controls}
The covariance matrix $\Sigma_C$ is positive semidefinite, $\kappa_C$ lies in its range,
and all minimisers of $\mathcal J$ are $\Sigma_C^\dagger \kappa_C+\ker \Sigma_C$, where $\Sigma_C^\dagger$
is the Moore--Penrose inverse. For every fixed coefficient vector,
\begin{equation}
 \mathcal J(\beta)=\mathcal J_{\mathcal G}(\beta)
 +\mathbb E_Q\!\left[
 \{Y-m-\beta^{\mathsf T}(C-C_{\mathcal G})\}^{2}\right].
 \label{eq:cgbm-joint-projection}
\end{equation}
Consequently $\min_\beta \mathcal J_{\mathcal G}(\beta)\leq\min_\beta \mathcal J(\beta)$.
If $\beta_{\rm raw}$ and $\beta_{\mathcal G}$ are population minimisers,
then
\begin{align}
 \mathcal J_{\mathcal G}(\beta_{\rm raw})-
 \mathcal J_{\mathcal G}(\beta_{\mathcal G})
 &=(\beta_{\rm raw}-\beta_{\mathcal G})^{\mathsf T}
 \Sigma_{C,\mathcal G}(\beta_{\rm raw}-\beta_{\mathcal G}),
 \label{eq:cgbm-coefficient-penalty}\\
 \Sigma_{C,\mathcal G}&=\mathbb E_Q[C_{\mathcal G}C_{\mathcal G}^{\mathsf T}].\notag
\end{align}
Thus projection is linear at fixed coefficients, but the optimal raw
coefficient need not remain optimal after conditioning.
\end{proposition}

For a terminal claim with horizon $T$ and initial spot $S_0=S$, set
$Y=e^{-rT}F(S_T)$ and choose $t=uT$ with $0<u<1$. Independent Brownian
increments give
\begin{align}
 m_t&=\mathbb E_Q[Y\mid\mathcal F_t]
 =\mathbb E_Q[Y\mid\sigma(S_t)]
 =e^{-rt}V(S_t,T-t),\label{eq:cgbm-conditional-price}\\
 C&=e^{-rT}S_T-S_0e^{-qT},\notag\\
 c_t&=\mathbb E_Q[C\mid\mathcal F_t]
 =e^{-rt-q(T-t)}S_t-S_0e^{-qT}.
 \label{eq:cgbm-conditional-control}
\end{align}
The Brownian filtrations $\mathcal F_t$ are nested; the individual
$\sigma(S_t)$ need not be. This distinction justifies comparing
population conditional variances across observation times. In particular,
\begin{equation}
 \operatorname{Var}_Q(c_t)
 =S_0^2e^{-2qT}(e^{\sigma^2t}-1).
 \label{eq:cgbm-conditional-control-variance}
\end{equation}
At $t=0$, $m_0=V(S_0,T)$ is the already available analytic price.
This zero-variance endpoint requires evaluating that price and is not a
new valuation algorithm.

I fit coefficients on a pilot independent of the production sample and
bound their norm before reuse. Appendix~\ref{app:cgbm-pilot} proves
conditional unbiasedness and finite variance for this procedure. The
population inequality in Proposition~\ref{prop:cgbm-projected-controls}
does not guarantee that a separately fitted conditional coefficient wins
on a finite production sample. The comparison therefore charges pilot
observations, terminal simulation and conditional-price evaluations
separately. The call-spread formula remains the strongest price baseline.
Unbiasedness here concerns exact projections; evaluating them in
floating-point arithmetic is not itself an unbiasedness certificate.

For the quantum basket method, the auxiliary-variance proposition already
uses the established total-variance identity. The present projection
argument explains that identity as a least-squares result and shows why
controls must be transformed jointly with the score. It does not change
the earlier finite-law or gate certificates. Reducing intrinsic variance
alone proves neither an overall computational advantage nor a quantum
advantage: the conditional evaluator and its required accuracy must also
be charged.
\subsection{Matched-information comparisons and evaluator work}

I reuse the seven synthetic contracts in Table~\ref{tab:cgbm-contracts}
and choose $u=0.25,0.5,0.9$. Every method has the same model parameters,
known stock-control means and analytic pricing function. The study uses
independent pilots of 8,192 observations and production samples of 65,536,
with coefficient norm bound $10^6$ and seed 840731. Methods share production
draws within each split. The split simulations have correct marginal
Brownian laws, but different observation times are not coupled into one
sampled Brownian path.

I compare raw payoffs, exact conditional valuation, raw pilot controls,
conditional valuation with the raw pilot coefficient, separately refitted
conditional controls, and the unbiased comparator that conditions the
payoff but leaves its control unprojected. The analytic call spread is
included in every case as the zero-sampling-error reference. There are
147 comparison rows, including that reference. Thirty-two additional
independent pilot fits per case/split give 672 coefficient-risk diagnostics.

Population raw moments use the truncated lognormal formulae derived in
Appendix~\ref{app:cgbm-moment-evaluation}. Conditional moments use
one-dimensional Gaussian quadrature at orders 160 and 320. These numerical
integrals are diagnostics, not rigorous interval enclosures. The maximum
refinement difference across recorded smooth moments is approximately
$3.83\times10^{-8}$; a different representation integrates the analytic
residual-life variance to check the projection decomposition.
Its maximum discrepancy is approximately $1.21\times10^{-10}$.
Coarse $80\times80$ quadrature of the kinked joint payoff instead has an
identity discrepancy as large as $4.62$ and remains an explicitly failed
resolution diagnostic. It is not used as validation of the smooth results.

\begin{table}[tb]\centering\small
\begin{tabular}{lrrrr}\toprule
Contract & Raw optimal & Conditional optimal & Ratio & Unprojected penalty\\\midrule
Lower kink & $14.303$ & $1.1166$ & 12.810 & $24.327$ \\
Upper kink & $21.919$ & $2.9954$ & 7.318 & $21.04$ \\
Inside spread & $17.683$ & $1.5105$ & 11.707 & $27.345$ \\
Deep cap & $7.0584$ & $1.1945$ & 5.909 & $0.71552$ \\
Short dated & $0.16239$ & $0.036354$ & 4.467 & $0.25742$ \\
Negative rate/dividend & $17.781$ & $1.2724$ & 13.975 & $25.975$ \\
High volatility & $70.98$ & $5.5678$ & 12.748 & $44.194$ \\
\bottomrule\end{tabular}
\caption{Population-variance diagnostics at $t=T/2$, in squared monetary units. Conditional quantities are approximated by order-320 Gaussian quadrature. The ratio is optimal raw-control variance divided by optimal conditional-control variance. The final column is the additional variance from leaving the stock control unprojected at the raw population coefficient. These are oracle comparisons, not finite-pilot efficiency guarantees.}
\label{tab:cgbm-projection-variance}\end{table}

Across the 21 interior cases/splits, the numerical oracle variance ratios
range from 1.2866 to 93.7667. They illustrate the population inequality,
without proving a universal numerical gain of that magnitude. The
conditional pilot fit has greater computed population excess variance than the
raw pilot coefficient reused on the projected control in 152 of 672
independent pilot realisations. This is a descriptive diagnostic of
coefficient-estimation risk, not a general probability or efficiency claim.
The fixed-coefficient projection inequality and the oracle optimum
inequality remain separate from this fitted-coefficient comparison.

\begin{figure}[tb]
\centering
\setlength{\unitlength}{1mm}
\begin{picture}(160,66)%
{\color[HTML]{E5E5E5}\put(20.000,21){\line(0,1){36}}}\put(20.000,16){\makebox(0,3){\scriptsize 0}}%
{\color[HTML]{E5E5E5}\put(20,21.000){\line(1,0){115}}}\put(17,19.500){\makebox(0,3)[r]{\scriptsize 0}}%
{\color[HTML]{E5E5E5}\put(48.750,21){\line(0,1){36}}}\put(48.750,16){\makebox(0,3){\scriptsize 0.25}}%
{\color[HTML]{E5E5E5}\put(20,30.000){\line(1,0){115}}}\put(17,28.500){\makebox(0,3)[r]{\scriptsize 0.25}}%
{\color[HTML]{E5E5E5}\put(77.500,21){\line(0,1){36}}}\put(77.500,16){\makebox(0,3){\scriptsize 0.5}}%
{\color[HTML]{E5E5E5}\put(20,39.000){\line(1,0){115}}}\put(17,37.500){\makebox(0,3)[r]{\scriptsize 0.5}}%
{\color[HTML]{E5E5E5}\put(106.250,21){\line(0,1){36}}}\put(106.250,16){\makebox(0,3){\scriptsize 0.75}}%
{\color[HTML]{E5E5E5}\put(20,48.000){\line(1,0){115}}}\put(17,46.500){\makebox(0,3)[r]{\scriptsize 0.75}}%
{\color[HTML]{E5E5E5}\put(135.000,21){\line(0,1){36}}}\put(135.000,16){\makebox(0,3){\scriptsize 1}}%
{\color[HTML]{E5E5E5}\put(20,57.000){\line(1,0){115}}}\put(17,55.500){\makebox(0,3)[r]{\scriptsize 1}}%
\put(20,60){\makebox(115,4){\small Retained variance $\operatorname{Var}(m_t)/\operatorname{Var}(Y)$}}%
\put(20,12){\makebox(115,3){\scriptsize Observation fraction $t/T$}}%
{\color[HTML]{0072B2}\linethickness{0.5pt}\qbezier(20.000,21.000)(34.375,24.383)(48.750,27.766)}%
{\color[HTML]{0072B2}\linethickness{0.5pt}\qbezier(48.750,27.766)(63.125,31.474)(77.500,35.183)}%
{\color[HTML]{0072B2}\linethickness{0.5pt}\qbezier(77.500,35.183)(100.500,42.974)(123.500,50.766)}%
{\color[HTML]{0072B2}\linethickness{0.5pt}\qbezier(123.500,50.766)(129.250,53.883)(135.000,57.000)}%
{\color[HTML]{0072B2}\put(20.000,21.000){\circle*{.8}}}%
{\color[HTML]{0072B2}\put(48.750,27.766){\circle*{.8}}}%
{\color[HTML]{0072B2}\put(77.500,35.183){\circle*{.8}}}%
{\color[HTML]{0072B2}\put(123.500,50.766){\circle*{.8}}}%
{\color[HTML]{0072B2}\put(135.000,57.000){\circle*{.8}}}%
{\color[HTML]{0072B2}\put(20,7){\line(1,0){5}}}\put(26,5.500){\makebox(0,3)[l]{\scriptsize Lower kink}}%
{\color[HTML]{D55E00}\linethickness{0.5pt}\qbezier(20.000,21.000)(34.375,22.570)(48.750,24.141)}%
{\color[HTML]{D55E00}\linethickness{0.5pt}\qbezier(48.750,24.141)(63.125,26.921)(77.500,29.701)}%
{\color[HTML]{D55E00}\linethickness{0.5pt}\qbezier(77.500,29.701)(100.500,38.747)(123.500,47.793)}%
{\color[HTML]{D55E00}\linethickness{0.5pt}\qbezier(123.500,47.793)(129.250,52.397)(135.000,57.000)}%
{\color[HTML]{D55E00}\put(20.000,21.000){\circle*{.8}}}%
{\color[HTML]{D55E00}\put(48.750,24.141){\circle*{.8}}}%
{\color[HTML]{D55E00}\put(77.500,29.701){\circle*{.8}}}%
{\color[HTML]{D55E00}\put(123.500,47.793){\circle*{.8}}}%
{\color[HTML]{D55E00}\put(135.000,57.000){\circle*{.8}}}%
{\color[HTML]{D55E00}\put(80,7){\line(1,0){5}}}\put(86,5.500){\makebox(0,3)[l]{\scriptsize Deep cap}}%
{\color[HTML]{009E73}\linethickness{0.5pt}\qbezier(20.000,21.000)(34.375,24.633)(48.750,28.265)}%
{\color[HTML]{009E73}\linethickness{0.5pt}\qbezier(48.750,28.265)(63.125,32.406)(77.500,36.548)}%
{\color[HTML]{009E73}\linethickness{0.5pt}\qbezier(77.500,36.548)(100.500,44.432)(123.500,52.316)}%
{\color[HTML]{009E73}\linethickness{0.5pt}\qbezier(123.500,52.316)(129.250,54.658)(135.000,57.000)}%
{\color[HTML]{009E73}\put(20.000,21.000){\circle*{.8}}}%
{\color[HTML]{009E73}\put(48.750,28.265){\circle*{.8}}}%
{\color[HTML]{009E73}\put(77.500,36.548){\circle*{.8}}}%
{\color[HTML]{009E73}\put(123.500,52.316){\circle*{.8}}}%
{\color[HTML]{009E73}\put(135.000,57.000){\circle*{.8}}}%
{\color[HTML]{009E73}\put(20,3){\line(1,0){5}}}\put(26,1.500){\makebox(0,3)[l]{\scriptsize Short dated}}%
{\color[HTML]{CC79A7}\linethickness{0.5pt}\qbezier(20.000,21.000)(34.375,23.829)(48.750,26.659)}%
{\color[HTML]{CC79A7}\linethickness{0.5pt}\qbezier(48.750,26.659)(63.125,29.884)(77.500,33.109)}%
{\color[HTML]{CC79A7}\linethickness{0.5pt}\qbezier(77.500,33.109)(100.500,40.544)(123.500,47.978)}%
{\color[HTML]{CC79A7}\linethickness{0.5pt}\qbezier(123.500,47.978)(129.250,52.489)(135.000,57.000)}%
{\color[HTML]{CC79A7}\put(20.000,21.000){\circle*{.8}}}%
{\color[HTML]{CC79A7}\put(48.750,26.659){\circle*{.8}}}%
{\color[HTML]{CC79A7}\put(77.500,33.109){\circle*{.8}}}%
{\color[HTML]{CC79A7}\put(123.500,47.978){\circle*{.8}}}%
{\color[HTML]{CC79A7}\put(135.000,57.000){\circle*{.8}}}%
{\color[HTML]{CC79A7}\put(80,3){\line(1,0){5}}}\put(86,1.500){\makebox(0,3)[l]{\scriptsize High volatility}}%
\end{picture}
\caption{Retained conditional-price variance for selected synthetic contracts. Interior markers use the one-dimensional moment diagnostics; endpoints follow the analytic model. Lines join evaluated points. Increasing the observation time retains more terminal-price randomness: estimator variance rises while payoff-prediction loss falls. These are distinct quantities. All seven contracts remain in the comparison files; the figure reports no computational speed or quantum advantage.}
\label{fig:cgbm-conditioning}
\end{figure}
A stand-alone raw estimate requires one Gaussian variate, one terminal
stock exponential and one capped payoff evaluation per observation. A
conditional estimate requires one Gaussian variate, one intermediate stock
exponential and four normal-CDF evaluations. Its independent conditional
pilot has the same conditional-evaluation requirement. Reusing a raw pilot
coefficient charges the raw pilot, while the unprojected-control comparator
also needs the additional terminal simulation. Per-policy counts and the
actual shared diagnostic generation counts are retained separately; shared
code execution is not presented as individual estimator timings.

For a population variance $v$ and per-observation work $c$, an idealised
large-sample work proxy at target mean-square sampling error $\epsilon^2$
is $cv/\epsilon^2$. The numerical variance ratio can therefore identify a
break-even production-cost ratio, but it excludes pilots, finite confidence,
kernel error and arithmetic. No CDF-to-payoff work conversion or wall-clock
superiority is established. The available analytic price still dominates
sampling in this model.

The conditional call-spread evaluator uses survival-function differences
for intervals in the upper normal tail and avoids subtracting two large
stock-valued calls. A range fallback removes tiny floating-point violations:
49 production states, four main-pilot states and 149 repeated-pilot states
use it, with largest recorded discounted adjustment
$3.46\times10^{-15}$. Those observations do not bound the complete kernel
error or prove numerical unbiasedness. The oracle/projection statements
refer to the exact conditional operator; the reported arithmetic and
quadrature diagnostics retain their narrower scope.

\section{Physical estimation and volatility uncertainty}
\label{sec:cgbm-estimation}

I now distinguish a stipulated pricing parameter from an estimated one.
The estimation experiment uses synthetic ex-dividend observations;
it tests statistical reasoning and the propagation of parameter uncertainty,
without establishing empirical model adequacy. I retain constant coefficients,
known dividends, equally spaced noise-free observations and independent
Gaussian increments. These assumptions exclude stochastic volatility,
serial dependence and market microstructure noise.

For $d$ assets under $P$, write
$dS_j/S_j=(\mu_j-q_j)\,dt+(L\,dW^P)_j$, where
$\Sigma=LL^\top\succ0$ and $\mu_j$ is a total-return drift. For spacing
$\delta>0$, the $n\geq2$ observed log-return vectors satisfy
\begin{equation}
 R_i\sim N_d(a\delta,\Sigma\delta),\qquad
 a_j=\mu_j-q_j-\tfrac12\Sigma_{jj},\qquad
 C=\sum_{i=1}^n(R_i-\bar R)(R_i-\bar R)^\top.
 \label{eq:cgbm-statistical-model}
\end{equation}
Appendix~\ref{app:cgbm-estimation} derives the likelihood and proves
\begin{equation}
 \widehat a=\bar R/\delta,\qquad
 \widehat\Sigma_{\mathrm{ML}}=C/(n\delta),\qquad
 \widehat\Sigma_{\mathrm U}=C/((n-1)\delta).
 \label{eq:cgbm-covariance-estimators}
\end{equation}
The full covariance maximum exists uniquely precisely when $C\succ0$.
When $n\leq d$, its rank is at most $n-1<d$ and the likelihood is
unbounded over positive-definite covariance matrices. Positive-definite
shrinkage therefore changes the estimation rule; it does not recover
a nonexistent unconstrained maximum.

The normal-sample laws have established foundations \citep{student1908,wishart1928}.
Gaussian row rotations give independent sample means and scatters, and
the exact Wishart law. They yield marginal Student-$t$ intervals for
$a_j$ and chi-squared intervals for $\sigma_j=\sqrt{\Sigma_{jj}}$.
The physical total-return estimator
$\widehat\mu_{\mathrm U,j}=\widehat a_j+q_j+
\widehat\Sigma_{\mathrm U,jj}/2$ is unbiased;
using the maximum-likelihood covariance instead gives bias
$-\Sigma_{jj}/(2n)$. The Student-$t$ interval concerns $a_j$, not
this corrected total-return drift. At fixed observation horizon
$T_{\mathrm{obs}}=n\delta$, the variance of $\widehat a_j$ is
$\Sigma_{jj}/T_{\mathrm{obs}}$: more frequent ideal observations
improve covariance precision but do not remove drift uncertainty.
The sampling-frequency distinction also has a published finance antecedent
\citep{merton1980}; the final journal metadata and working-paper reading
are distinguished in the source ledger.
Under the stipulated constant-diffusion change of measure,
covariance is shared by $P$ and $Q$, while the pricing drift remains
$r-q_j$ regardless of the estimated physical drift.

\begin{proposition}[The image of a volatility confidence interval]
\label{prop:cgbm-volatility-image}
For fixed known $S,K,H,\tau,r,q$, set
\begin{equation}
 c_* =\frac{\log(K(K+H)/S^2)}{\tau}-2(r-q).
 \label{eq:cgbm-volatility-critical}
\end{equation}
If $c_*\leq0$, the capped price is strictly decreasing in $\sigma>0$.
If $c_*>0$, it increases up to $\sigma_* =\sqrt{c_*}$ and decreases
thereafter. For an interval $[\ell,u]\subset(0,\infty)$, its minimum
price is the smaller endpoint price; its maximum is the largest price
among the endpoints and $\sigma_*$ when $\sigma_*\in[\ell,u]$.
Consequently, mapping an exact marginal volatility confidence interval
through this full image gives a price interval with at least the same
coverage under the stipulated model.
\end{proposition}

The proof uses the sign of vega in \eqref{eq:cgbm-vega}; no monotonicity
assumption is imposed on a capped claim. The synthetic comparison includes
an interval straddling $\sigma_*$, for which endpoint-only propagation
misses the upper price. At the true interior volatility maximum, endpoint-only propagation
of the exact chi-squared interval has zero coverage, whereas the full image
preserves the volatility interval's coverage. This is a financial consequence of changing
convexity across the two strikes. The coverage statement is marginal,
conditional on known contractual inputs and the model assumptions;
floating-point quantiles and prices are not outward-certified bounds.
Shrinkage does not inherit the exact chi-squared pivot, and statistical
uncertainty remains distinct from PDE and Monte Carlo numerical errors.
These are established estimation principles and a checked application
of the call-spread sensitivity, rather than a claim of theoretical priority.

Conditional numerical price guarantees at the interval endpoints and included
volatility maximum compose with statistical coverage to give at least
$1-\alpha-\beta$ overall coverage (Appendix~\ref{app:cgbm-estimation}).
A classical/quantum comparison must share estimation information and charge
all required pricing calls. The earlier floating-point experiment remains
a diagnostic comparison rather than a validated certificate.

I implement a validated bridge from rounded observations to a model-price
interval (Appendix~\ref{app:cgbm-validated-calibration}). Exact rational
scatter calculations retain observation-rounding error; outward
chi-squared quantile brackets and price enclosures then replace
uncertified floating-point evaluations. The maximum is evaluated at a
clamped stationary volatility, including uncertain endpoint coincidences.
The resulting guarantee concerns the stipulated Gaussian constant-diffusion
model, with every other price input known. It separates parameter
uncertainty from numerical enclosure width. A returned certificate has a
joint failure probability at most $\alpha$; this is not conditional
coverage after selecting successful computations. Resource exhaustion
produces an explicit abstention rather than a false numerical guarantee.

A separate certificate study retains 13 adverse price-image cases,
12 chi-squared quantile brackets and 180 synthetic datasets, with 23
independent implementation tests. Every dataset returns a numerical
certificate, while true-price containment is 165/180. The SciPy full-image
comparator gives the same count. In the near-stationary case, full-image
propagation contains 54/60 targets and endpoint-only propagation none.
These are diagnostics from approximate Gaussian generation and six-decimal
observation rounding, not an empirical 95 per cent coverage guarantee.
The gain is verified outward arithmetic, rather than increased observed
coverage. Inputs, failures, exact bounds and informed implementation
corrections are retained with the study.


\subsection{Synthetic estimation and price-coverage evidence}

I generate three-dimensional log returns using physical total drifts
$(0.06,0.04,0.08)$, dividend yields $(q,0.015,0.02)$ and volatilities
$(\sigma,0.8\sigma,1.2\sigma)$. Here $q,\sigma$ are the first-asset
values of each designed contract. The correlation matrix is
\begin{equation}
 \begin{pmatrix}1&0.4&-0.2\\0.4&1&0.25\\-0.2&0.25&1\end{pmatrix}.
 \label{eq:cgbm-synthetic-correlation}
\end{equation}
I retain the seven original contracts and add a first-asset contract with
$S=K=100$, $H=20$, $\tau=1$, $r=0.03$, $q=0.01$ and
$\sigma=\sqrt{\log(1.2)-0.04}=\sigma_*$.
Each case uses 1,000 independent synthetic datasets of 252 observations
at spacing $1/252$. Marginal intervals use $\alpha=0.05$; simultaneous
mean/diagonal-variance rectangles allocate $0.05/(2d)$ per interval.
The seed is 926341, with fixed case offsets. Code and summaries are
retained in \texttt{results/capped\_gbm/calibration}; raw returns are not saved.

\begin{table}[tb]\centering\small
\begin{tabular}{lrrr}\toprule
Contract & Full price image & Endpoints only & Plug-in price RMSE\\\midrule
Lower kink & 937/1000 & 937/1000 & 0.0760989 \\
Upper kink & 946/1000 & 946/1000 & 0.189076 \\
Inside spread & 954/1000 & 954/1000 & 0.0676302 \\
Deep cap & 951/1000 & 951/1000 & 0.123294 \\
Short dated & 959/1000 & 959/1000 & 0.0231021 \\
Negative rate/dividend & 965/1000 & 965/1000 & 0.08144 \\
High volatility & 939/1000 & 939/1000 & 0.129354 \\
Volatility maximum & 948/1000 & 0/1000 & 0.00447807 \\
\bottomrule\end{tabular}
\caption{Observed coverage and plug-in error across independent synthetic datasets. Coverage counts concern marginal 95 per cent volatility intervals and their price images. RMSE is in monetary units. Observed counts are numerical diagnostics, not statements that each experiment attained at least 95 per cent coverage.}
\label{tab:cgbm-calibration-coverage}\end{table}

For the volatility-maximum contract, the full image covers 948 of 1,000
datasets; endpoint-only propagation covers none. The pointwise 95 per cent
Clopper--Pearson interval for the former coverage frequency is
$[0.93236,0.96092]$, and the zero-success interval has upper endpoint
$0.003682$ \citep{clopperpearson1934}. Across all cases the corrected
price counts range from 937 to 965; the simultaneous parameter-rectangle
counts range from 944 to 962. These fluctuations remain visible. Neither
the repeated-data binomial intervals nor the individual parameter intervals
are claimed to be simultaneously calibrated across every reported diagnostic.
The exact probability proof and numerical coverage check are distinct.

\begin{figure}[tb]
\centering
\setlength{\unitlength}{1mm}
\begin{picture}(160,66)%
{\color[HTML]{DDDDDD}\put(20.000,19){\line(0,1){36}}}\put(20.000,15){\makebox(0,3){\scriptsize 0.3395}}%
{\color[HTML]{DDDDDD}\put(77.500,19){\line(0,1){36}}}\put(77.500,15){\makebox(0,3){\scriptsize 0.3773}}%
{\color[HTML]{DDDDDD}\put(135.000,19){\line(0,1){36}}}\put(135.000,15){\makebox(0,3){\scriptsize 0.4150}}%
{\color[HTML]{DDDDDD}\put(20,19.000){\line(1,0){115}}}\put(18,17.500){\makebox(0,3)[r]{\scriptsize 6.808}}%
{\color[HTML]{DDDDDD}\put(20,28.000){\line(1,0){115}}}\put(18,26.500){\makebox(0,3)[r]{\scriptsize 6.813}}%
{\color[HTML]{DDDDDD}\put(20,37.000){\line(1,0){115}}}\put(18,35.500){\makebox(0,3)[r]{\scriptsize 6.818}}%
{\color[HTML]{DDDDDD}\put(20,46.000){\line(1,0){115}}}\put(18,44.500){\makebox(0,3)[r]{\scriptsize 6.822}}%
{\color[HTML]{DDDDDD}\put(20,55.000){\line(1,0){115}}}\put(18,53.500){\makebox(0,3)[r]{\scriptsize 6.827}}%
{\color[HTML]{0072B2}\linethickness{0.6pt}\qbezier(20.000,22.874)(20.719,23.606)(21.438,24.339)}%
{\color[HTML]{0072B2}\linethickness{0.6pt}\qbezier(21.438,24.339)(22.156,25.050)(22.875,25.762)}%
{\color[HTML]{0072B2}\linethickness{0.6pt}\qbezier(22.875,25.762)(23.594,26.452)(24.313,27.142)}%
{\color[HTML]{0072B2}\linethickness{0.6pt}\qbezier(24.313,27.142)(25.031,27.811)(25.750,28.480)}%
{\color[HTML]{0072B2}\linethickness{0.6pt}\qbezier(25.750,28.480)(26.469,29.128)(27.188,29.776)}%
{\color[HTML]{0072B2}\linethickness{0.6pt}\qbezier(27.188,29.776)(27.906,30.404)(28.625,31.031)}%
{\color[HTML]{0072B2}\linethickness{0.6pt}\qbezier(28.625,31.031)(29.344,31.638)(30.062,32.245)}%
{\color[HTML]{0072B2}\linethickness{0.6pt}\qbezier(30.062,32.245)(30.781,32.831)(31.500,33.418)}%
{\color[HTML]{0072B2}\linethickness{0.6pt}\qbezier(31.500,33.418)(32.219,33.984)(32.937,34.550)}%
{\color[HTML]{0072B2}\linethickness{0.6pt}\qbezier(32.937,34.550)(33.656,35.096)(34.375,35.642)}%
{\color[HTML]{0072B2}\linethickness{0.6pt}\qbezier(34.375,35.642)(35.094,36.168)(35.812,36.694)}%
{\color[HTML]{0072B2}\linethickness{0.6pt}\qbezier(35.812,36.694)(36.531,37.201)(37.250,37.707)}%
{\color[HTML]{0072B2}\linethickness{0.6pt}\qbezier(37.250,37.707)(37.969,38.193)(38.688,38.680)}%
{\color[HTML]{0072B2}\linethickness{0.6pt}\qbezier(38.688,38.680)(39.406,39.147)(40.125,39.614)}%
{\color[HTML]{0072B2}\linethickness{0.6pt}\qbezier(40.125,39.614)(40.844,40.062)(41.563,40.510)}%
{\color[HTML]{0072B2}\linethickness{0.6pt}\qbezier(41.563,40.510)(42.281,40.938)(43.000,41.367)}%
{\color[HTML]{0072B2}\linethickness{0.6pt}\qbezier(43.000,41.367)(43.719,41.777)(44.437,42.186)}%
{\color[HTML]{0072B2}\linethickness{0.6pt}\qbezier(44.437,42.186)(45.156,42.577)(45.875,42.967)}%
{\color[HTML]{0072B2}\linethickness{0.6pt}\qbezier(45.875,42.967)(46.594,43.339)(47.312,43.711)}%
{\color[HTML]{0072B2}\linethickness{0.6pt}\qbezier(47.312,43.711)(48.031,44.064)(48.750,44.417)}%
{\color[HTML]{0072B2}\linethickness{0.6pt}\qbezier(48.750,44.417)(49.469,44.752)(50.187,45.086)}%
{\color[HTML]{0072B2}\linethickness{0.6pt}\qbezier(50.187,45.086)(50.906,45.403)(51.625,45.719)}%
{\color[HTML]{0072B2}\linethickness{0.6pt}\qbezier(51.625,45.719)(52.344,46.017)(53.062,46.315)}%
{\color[HTML]{0072B2}\linethickness{0.6pt}\qbezier(53.062,46.315)(53.781,46.595)(54.500,46.875)}%
{\color[HTML]{0072B2}\linethickness{0.6pt}\qbezier(54.500,46.875)(55.219,47.137)(55.938,47.399)}%
{\color[HTML]{0072B2}\linethickness{0.6pt}\qbezier(55.938,47.399)(56.656,47.644)(57.375,47.888)}%
{\color[HTML]{0072B2}\linethickness{0.6pt}\qbezier(57.375,47.888)(58.094,48.114)(58.813,48.341)}%
{\color[HTML]{0072B2}\linethickness{0.6pt}\qbezier(58.813,48.341)(59.531,48.550)(60.250,48.759)}%
{\color[HTML]{0072B2}\linethickness{0.6pt}\qbezier(60.250,48.759)(60.969,48.950)(61.687,49.142)}%
{\color[HTML]{0072B2}\linethickness{0.6pt}\qbezier(61.687,49.142)(62.406,49.316)(63.125,49.490)}%
{\color[HTML]{0072B2}\linethickness{0.6pt}\qbezier(63.125,49.490)(63.844,49.647)(64.562,49.804)}%
{\color[HTML]{0072B2}\linethickness{0.6pt}\qbezier(64.562,49.804)(65.281,49.944)(66.000,50.084)}%
{\color[HTML]{0072B2}\linethickness{0.6pt}\qbezier(66.000,50.084)(66.719,50.208)(67.438,50.331)}%
{\color[HTML]{0072B2}\linethickness{0.6pt}\qbezier(67.438,50.331)(68.156,50.437)(68.875,50.543)}%
{\color[HTML]{0072B2}\linethickness{0.6pt}\qbezier(68.875,50.543)(69.594,50.633)(70.312,50.722)}%
{\color[HTML]{0072B2}\linethickness{0.6pt}\qbezier(70.312,50.722)(71.031,50.795)(71.750,50.869)}%
{\color[HTML]{0072B2}\linethickness{0.6pt}\qbezier(71.750,50.869)(72.469,50.925)(73.188,50.982)}%
{\color[HTML]{0072B2}\linethickness{0.6pt}\qbezier(73.188,50.982)(73.906,51.022)(74.625,51.062)}%
{\color[HTML]{0072B2}\linethickness{0.6pt}\qbezier(74.625,51.062)(75.344,51.086)(76.063,51.110)}%
{\color[HTML]{0072B2}\linethickness{0.6pt}\qbezier(76.063,51.110)(76.781,51.118)(77.500,51.126)}%
{\color[HTML]{0072B2}\linethickness{0.6pt}\qbezier(77.500,51.126)(78.219,51.119)(78.937,51.111)}%
{\color[HTML]{0072B2}\linethickness{0.6pt}\qbezier(78.937,51.111)(79.656,51.087)(80.375,51.063)}%
{\color[HTML]{0072B2}\linethickness{0.6pt}\qbezier(80.375,51.063)(81.094,51.023)(81.812,50.984)}%
{\color[HTML]{0072B2}\linethickness{0.6pt}\qbezier(81.812,50.984)(82.531,50.929)(83.250,50.873)}%
{\color[HTML]{0072B2}\linethickness{0.6pt}\qbezier(83.250,50.873)(83.969,50.803)(84.687,50.732)}%
{\color[HTML]{0072B2}\linethickness{0.6pt}\qbezier(84.687,50.732)(85.406,50.646)(86.125,50.560)}%
{\color[HTML]{0072B2}\linethickness{0.6pt}\qbezier(86.125,50.560)(86.844,50.458)(87.562,50.357)}%
{\color[HTML]{0072B2}\linethickness{0.6pt}\qbezier(87.562,50.357)(88.281,50.240)(89.000,50.124)}%
{\color[HTML]{0072B2}\linethickness{0.6pt}\qbezier(89.000,50.124)(89.719,49.992)(90.438,49.860)}%
{\color[HTML]{0072B2}\linethickness{0.6pt}\qbezier(90.438,49.860)(91.156,49.714)(91.875,49.567)}%
{\color[HTML]{0072B2}\linethickness{0.6pt}\qbezier(91.875,49.567)(92.594,49.406)(93.313,49.244)}%
{\color[HTML]{0072B2}\linethickness{0.6pt}\qbezier(93.313,49.244)(94.031,49.068)(94.750,48.891)}%
{\color[HTML]{0072B2}\linethickness{0.6pt}\qbezier(94.750,48.891)(95.469,48.701)(96.187,48.510)}%
{\color[HTML]{0072B2}\linethickness{0.6pt}\qbezier(96.187,48.510)(96.906,48.304)(97.625,48.099)}%
{\color[HTML]{0072B2}\linethickness{0.6pt}\qbezier(97.625,48.099)(98.344,47.879)(99.062,47.659)}%
{\color[HTML]{0072B2}\linethickness{0.6pt}\qbezier(99.062,47.659)(99.781,47.424)(100.500,47.190)}%
{\color[HTML]{0072B2}\linethickness{0.6pt}\qbezier(100.500,47.190)(101.219,46.941)(101.938,46.693)}%
{\color[HTML]{0072B2}\linethickness{0.6pt}\qbezier(101.938,46.693)(102.656,46.430)(103.375,46.167)}%
{\color[HTML]{0072B2}\linethickness{0.6pt}\qbezier(103.375,46.167)(104.094,45.891)(104.813,45.614)}%
{\color[HTML]{0072B2}\linethickness{0.6pt}\qbezier(104.813,45.614)(105.531,45.323)(106.250,45.032)}%
{\color[HTML]{0072B2}\linethickness{0.6pt}\qbezier(106.250,45.032)(106.969,44.728)(107.688,44.423)}%
{\color[HTML]{0072B2}\linethickness{0.6pt}\qbezier(107.688,44.423)(108.406,44.105)(109.125,43.786)}%
{\color[HTML]{0072B2}\linethickness{0.6pt}\qbezier(109.125,43.786)(109.844,43.454)(110.562,43.122)}%
{\color[HTML]{0072B2}\linethickness{0.6pt}\qbezier(110.562,43.122)(111.281,42.777)(112.000,42.431)}%
{\color[HTML]{0072B2}\linethickness{0.6pt}\qbezier(112.000,42.431)(112.719,42.072)(113.437,41.713)}%
{\color[HTML]{0072B2}\linethickness{0.6pt}\qbezier(113.437,41.713)(114.156,41.340)(114.875,40.968)}%
{\color[HTML]{0072B2}\linethickness{0.6pt}\qbezier(114.875,40.968)(115.594,40.582)(116.312,40.196)}%
{\color[HTML]{0072B2}\linethickness{0.6pt}\qbezier(116.312,40.196)(117.031,39.797)(117.750,39.398)}%
{\color[HTML]{0072B2}\linethickness{0.6pt}\qbezier(117.750,39.398)(118.469,38.986)(119.188,38.574)}%
{\color[HTML]{0072B2}\linethickness{0.6pt}\qbezier(119.188,38.574)(119.906,38.149)(120.625,37.724)}%
{\color[HTML]{0072B2}\linethickness{0.6pt}\qbezier(120.625,37.724)(121.344,37.286)(122.063,36.847)}%
{\color[HTML]{0072B2}\linethickness{0.6pt}\qbezier(122.063,36.847)(122.781,36.396)(123.500,35.946)}%
{\color[HTML]{0072B2}\linethickness{0.6pt}\qbezier(123.500,35.946)(124.219,35.482)(124.938,35.018)}%
{\color[HTML]{0072B2}\linethickness{0.6pt}\qbezier(124.938,35.018)(125.656,34.542)(126.375,34.065)}%
{\color[HTML]{0072B2}\linethickness{0.6pt}\qbezier(126.375,34.065)(127.094,33.576)(127.813,33.087)}%
{\color[HTML]{0072B2}\linethickness{0.6pt}\qbezier(127.813,33.087)(128.531,32.586)(129.250,32.084)}%
{\color[HTML]{0072B2}\linethickness{0.6pt}\qbezier(129.250,32.084)(129.969,31.570)(130.688,31.056)}%
{\color[HTML]{0072B2}\linethickness{0.6pt}\qbezier(130.688,31.056)(131.406,30.530)(132.125,30.004)}%
{\color[HTML]{0072B2}\linethickness{0.6pt}\qbezier(132.125,30.004)(132.844,29.465)(133.562,28.927)}%
{\color[HTML]{0072B2}\linethickness{0.6pt}\qbezier(133.562,28.927)(134.281,28.376)(135.000,27.825)}%
{\color[HTML]{D55E00}\linethickness{0.6pt}\put(20,27.825){\line(1,0){115}}}%
{\color[HTML]{009E73}\put(77.500,51.126){\circle*{1.2}}}%
\put(4,37){\makebox(0,0){\rotatebox{90}{\scriptsize Model price}}}%
\put(20,60){\makebox(115,4){\small A deterministic volatility-image counterexample}}%
\put(20,11){\makebox(115,3){\scriptsize Volatility $\sigma$}}%
{\color[HTML]{0072B2}\put(20,7){\line(1,0){5}}}\put(26,5.5){\makebox(0,3)[l]{\scriptsize Capped price}}%
{\color[HTML]{D55E00}\put(80,7){\line(1,0){5}}}\put(86,5.5){\makebox(0,3)[l]{\scriptsize Endpoint-only upper}}%
{\color[HTML]{009E73}\put(20,3){\circle*{.8}}}\put(26,1.5){\makebox(0,3)[l]{\scriptsize Interior maximum}}%
\end{picture}
\caption{A deterministic interval $[0.9\sigma_*,1.1\sigma_*]$ for the synthetic volatility-maximum contract. No confidence level is assigned to this illustration. Its maximum price is attained inside the interval, so the endpoint-only upper value misses a legitimate price. The repeated-data confidence experiment uses the chi-squared intervals separately.}
\label{fig:cgbm-volatility-image}
\end{figure}
The deterministic counterexample in Figure~\ref{fig:cgbm-volatility-image}
uses $[0.9\sigma_*,1.1\sigma_*]$. Its largest endpoint price is
$6.81278566$, whereas the true interior maximum is $6.82481668$.
This gap is a missing extremum rather than Monte Carlo sampling error.
No statistical confidence level is assigned to that illustrative interval.

\begin{table}[tb]\centering\small
\begin{tabular}{rrlrr}\toprule
$n$ & Horizon & Spacing & Empirical log-drift SD & Theoretical SD\\\midrule
63 & 1 & $1/63$ & 0.20042 & 0.20000 \\
252 & 1 & $1/252$ & 0.19918 & 0.20000 \\
1008 & 1 & $1/1008$ & 0.20298 & 0.20000 \\
252 & 1 & $1/252$ & 0.20573 & 0.20000 \\
1008 & 4 & $1/252$ & 0.10021 & 0.10000 \\
\bottomrule\end{tabular}
\caption{First-asset log-drift uncertainty over 1,000 independent datasets per regime. The repeated one-year daily row belongs to the separate horizon comparison. Increasing frequency at a fixed one-year horizon does not reduce the drift floor; extending the observation horizon does. The result assumes ideal independent diffusion increments.}
\label{tab:cgbm-frequency-horizon}\end{table}

The experiment also keeps a singular three-observation/three-asset sample.
Its numerical scatter rank is two, and the unrestricted covariance maximum
does not exist. Trace-target shrinkage with intensity $0.1$ gives positive
eigenvalues, but changes the estimator and supplies no transplanted exact
interval law. Numerical rank uses floating-point tolerances; the theorem
concerns exact rank. The study records estimation bias, covariance error,
pivot moments and all interval failures. It estimates a stipulated synthetic
model, without validating market prices or identifying empirical volatility
risk premia. Pricing uses $r-q$, not a fitted physical return mean.

\section{Interpretation and reproducibility}

This companion closes a specific preparation gap by connecting an
admissible financial strategy, a PDE, a nonsmooth numerical convergence
proof and independently checked sensitivities. The bounded mixed-gamma
derivation gives a useful estimator and a finite-sample uncertainty
statement for this claim. It is an explicit specialisation of sensitivity
principles, not a claimed new quantum theory or a completed priority claim.
The comparison retains negative gamma/vega, non-monotone errors and
adverse control rankings where they occur.

The numerical pricing comparisons use a stipulated one-asset model with
known parameters. The separate estimation experiment fits synthetic
observations and propagates parameter uncertainty. It does not validate
market prices, solve a general portfolio/control problem or establish a
quantum advantage.
It also does not extend a one-asset analytic benchmark into a closed-form
solution for arbitrary baskets. Those boundaries are essential when
using the companion as evidence of financial mathematics and computation.

The engine is \texttt{capped\_gbm.py};
\texttt{run\_capped\_gbm\_study.py} reproduces the comparisons.
The tests include active-region Gaussian quadrature of the mixed gamma,
analytic sensitivity checks, the direct/sine-basis backward-Euler comparison
and projection/singular-control checks. Gaussian inference and price-image
checks use \path{gbm_calibration.py} and
\path{run_calibration_study.py}; their results and source ledger are
retained in \path{results/capped_gbm/calibration}. The conditioning
code and source ledger are retained in \texttt{results/capped\_gbm/conditioning}. The source ledgers record publication details, the versions actually
consulted and the precise reading scope. Detailed arguments follow in
the appendices. I retain this document as a completed companion study;
the main quantum basket paper remains a working paper.

\appendix


\section{Financial and analytic derivations}
\label{app:cgbm-financial}

\subsection{Change of measure and the self-financing equation}

For a concrete construction, suppose under a physical measure $P$ that
$dS_t=(\mu-q)S_t\,dt+\sigma S_t\,dW_t^P$ with constant total-return drift
$\mu$. Set $\vartheta=(\mu-r)/\sigma$. Novikov's condition follows from
$\exp(\vartheta^2T/2)<\infty$ on every finite horizon. Thus
\begin{equation}
 \frac{dQ}{dP}\bigg|_{\mathcal F_T}
 =\exp\{-\vartheta W_T^P-\vartheta^2T/2\},\qquad
 W_t^Q=W_t^P+\vartheta t
 \label{eq:cgbm-rn}
\end{equation}
defines an equivalent probability measure and yields
\eqref{eq:cgbm-law}. The stochastic exponential is strictly positive
and has expectation one by the Gaussian moment-generating function.
For this constant coefficient, the Brownian conclusion can also be
checked directly. On any deterministic partition, the density is the
product of factors
$\exp\{-\vartheta\Delta W^P-\vartheta^2\Delta t/2\}$.
Multiplying the independent $N(0,\Delta t)$ increment densities by these
factors and completing the square changes each mean to
$-\vartheta\Delta t$, leaves its variance $\Delta t$ and preserves their
factorisation. Thus the increments of $W^Q=W^P+\vartheta t$ are independent
centred normals with the Brownian variances. The same calculation
conditional on the preceding Brownian filtration proves independence
from its past. Continuity is inherited from $W^P$. Hence $W^Q$ is a
Brownian motion under $Q$, an elementary verification of this established
constant-coefficient measure change.
Consequently this construction does not depend on an unverified
calibration of $\mu$.

Let $B_t=e^{rt}$, let $\delta_t$ be the share holding and let $\gamma_t$
be the bank-account holding. A self-financing wealth $P_t=\delta_tS_t+
\gamma_tB_t$ receives asset dividends, so
\begin{equation}
 dP_t=\delta_t\,dS_t+q\delta_tS_t\,dt+
 r(P_t-\delta_tS_t)\,dt
 =rP_t\,dt+\delta_t\sigma S_t\,dW_t^Q.
 \label{eq:cgbm-self-financing}
\end{equation}
For an admissible strategy, the share holding is predictable and
\begin{equation}
 \int_0^T e^{-2rt}\delta_t^2\sigma^2S_t^2\,dt<\infty
 \quad Q\text{-almost surely},\label{eq:cgbm-admissible-integrand}
\end{equation}
while discounted wealth is bounded below by a deterministic constant. For such a strategy, discounted
wealth is a local martingale bounded below and hence a supermartingale.
One may see the latter implication by localisation and conditional
Fatou's lemma after adding the fixed lower bound. An arbitrage with zero
initial wealth and non-negative, sometimes strictly positive terminal
wealth would contradict the supermartingale expectation inequality.
This proves no arbitrage within this explicit admissible strategy class;
it does not assert a general fundamental theorem for arbitrary markets.

\subsection{The analytic call and replication}

The event $X>A$ is $Z>-d_{2,A}$. Completing the square gives
\begin{equation}
 \mathbb E[e^{vZ}\mathbf 1_{\{Z>z\}}]
 =e^{v^2/2}\Phi(v-z).
 \label{eq:cgbm-gaussian-tilt}
\end{equation}
Substitution in $D\mathbb E[(X-A)^+]$ yields
\eqref{eq:cgbm-call}; payoff linearity yields the call spread. The normal
identity
\begin{equation}
 S e^{-q\tau}\varphi(d_{1,A})=A D\varphi(d_{2,A})
 \label{eq:cgbm-density-identity}
\end{equation}
follows directly by subtracting the two squared normal arguments. It
cancels the derivative-of-$\Phi$ terms when differentiating the call
formula. The normal arguments have derivatives
\begin{align}
 \partial_S d_{1,A}&=\partial_S d_{2,A}=\frac1{Sv},\notag\\
 \partial_\sigma d_{1,A}&=\sqrt{\tau}-\frac{d_{1,A}}{\sigma},
 &\partial_\sigma d_{2,A}&=-\frac{d_{1,A}}{\sigma},\notag\\
 \partial_r d_{1,A}&=\partial_r d_{2,A}=\frac{\sqrt{\tau}}{\sigma},
 \label{eq:cgbm-d-derivatives}\\
 \partial_\tau d_{1,A}&=
 \frac{r-q+\sigma^2/2}{\sigma\sqrt{\tau}}
 -\frac{d_{1,A}}{2\tau},
 &\partial_\tau d_{2,A}&=\partial_\tau d_{1,A}
 -\frac{\sigma}{2\sqrt{\tau}}.
 \notag
\end{align}
In the spot derivative of $C_A$, the two density terms cancel exactly,
leaving $e^{-q\tau}\Phi(d_{1,A})$; another spot derivative gives
$e^{-q\tau}\varphi(d_{1,A})/(Sv)$. In the volatility derivative, their
combined contribution is
$S e^{-q\tau}\varphi(d_{1,A})
(\partial_\sigma d_{1,A}-\partial_\sigma d_{2,A})
=S e^{-q\tau}\varphi(d_{1,A})\sqrt{\tau}$.
The rate-argument derivatives cancel, leaving
$\tau A D\Phi(d_{2,A})$ from differentiating the discount factor.
Finally,
\begin{equation}
 \partial_\tau C_A=-q S e^{-q\tau}\Phi(d_{1,A})
 +r A D\Phi(d_{2,A})
 +\frac{S e^{-q\tau}\sigma\varphi(d_{1,A})}{2\sqrt{\tau}}.
 \label{eq:cgbm-call-time}
\end{equation}
Subtracting the corresponding strike-$J$ expressions from those at
strike $K$ proves every formula in
\eqref{eq:cgbm-delta}--\eqref{eq:cgbm-theta}.
The lognormal convolution is smooth at positive maturity, even though
$F$ has two kinks. Dominated convergence proves continuity at expiry:
$X\to S$ almost surely, $F$ is continuous and $0\leq F\leq H$.

To verify the heat transformation explicitly, first write
$x=\log(S/K)$ and $V=e^{-r\tau}u(\tau,x)$. The chain rule gives
$SV_S=e^{-r\tau}u_x$ and
$S^2V_{SS}=e^{-r\tau}(u_{xx}-u_x)$. Equation
\eqref{eq:cgbm-pde} therefore reduces to
$u_\tau=a u_{xx}+b u_x$. Setting $u(\tau,x)=w(\tau,x+b\tau)$ gives
$w_\tau=a w_{yy}$. Conversely, Gaussian convolution proves
\eqref{eq:cgbm-heat}; reversing the calculation proves
\eqref{eq:cgbm-pde} without assuming it in the valuation argument.

With calendar time $t$ and expiry $T$, apply It\^o's formula to
$V(S_t,T-t)$ on $[0,T-\varepsilon]$. The pricing equation gives
\begin{equation}
 dV(S_t,T-t)=rV(S_t,T-t)\,dt+
 \sigma S_tV_S(S_t,T-t)\,dW_t^Q.
 \label{eq:cgbm-replication-ito}
\end{equation}
This agrees with \eqref{eq:cgbm-self-financing} for $\delta=V_S$ and
$\gamma=(V-SV_S)/B$. Furthermore
$0\leq V_S\leq e^{-q(T-t)}$ by \eqref{eq:cgbm-delta}, and
the Gaussian moment-generating function gives the explicit second moment
\begin{equation}
 \mathbb E_Q S_t^2=S_0^2
 \exp\{[2(r-q)+\sigma^2]t\}.
 \label{eq:cgbm-asset-second-moment}
\end{equation}
It is finite and bounded on every finite horizon. Hence
$\mathbb E_Q\int_0^T e^{-2rt}\sigma^2 S_t^2 V_S^2\,dt<\infty$.
The discounted stochastic gain is therefore a true square-integrable
martingale. Continuity at expiry gives terminal wealth $F(S_T)$ as
$\varepsilon\downarrow0$. Its discounted wealth is non-negative, so the
strategy is admissible. Any bounded quoted claim valuation different from
the replication cost permits buying the cheaper position, selling the dearer
one and investing the initial difference; their terminal claim cash
flows cancel. This gives a strictly positive terminal gain with zero
initial wealth and a bounded-below discounted wealth process, contradicting
no arbitrage. Any other bounded classical solution with this terminal
payoff has a discounted martingale representation with the same
conditional terminal expectation; equivalently, apply the stopped
It\^o argument and bounded convergence. Its price equals
\eqref{eq:cgbm-price}. This proves the stated uniqueness.

\section{Numerical stability and convergence proofs}
\label{app:cgbm-numerical}

\subsection{The matrix maximum principle}

Let $A_h$ be the $(M-1)\times(M-1)$ matrix with diagonal
$1+2\lambda$ and off-diagonals $-\lambda$. If $A_h z\geq0$ componentwise
and $z$ had a negative minimum at an interior row, that row would satisfy
$(A_hz)_j\leq z_j<0$. At an edge row, the missing neighbour is interpreted
as zero and gives an even stronger contradiction. Thus $z\geq0$.
The same argument applied to a null vector and its negative proves that
$A_h$ is nonsingular; consequently $A_h^{-1}\geq0$. The row sums of
$A_h$ are one in interior rows and $1+\lambda$ at the two edge rows.
Therefore $A_h\mathbf 1\geq\mathbf 1$, whence
\begin{equation}
 0\leq A_h^{-1}\mathbf 1\leq\mathbf 1,
 \qquad \|A_h^{-1}\|_\infty\leq1.
 \label{eq:cgbm-mmatrix}
\end{equation}
This proves contraction for two solutions with equal boundary values.
For positivity and the upper bound, the actual right-hand side is the
previous interior vector plus $\lambda H$ at the last row. If the previous
vector lies between $0$ and $H\mathbf 1$, this right-hand side lies between
$0$ and $A_h(H\mathbf 1)$. Applying the non-negative inverse proves the
invariant interval by induction. No condition relating $k$ and $h$ is used.

The discrete eigenvectors are $\sin(j\ell\pi/M)$ and substitution in the
centred second difference gives $-\nu_\ell$. The affine lift
$\ell(y)=H(y-L)/(U-L)$ has zero discrete second difference. Each lifted
mode is consequently multiplied by $(1+ak\nu_\ell)^{-1}$ per step,
which proves \eqref{eq:cgbm-dst} after $N$ steps. This identity is in exact
arithmetic; the direct-step comparison checks implementation error.

\subsection{Mollification without a false smooth-payoff assumption}

Let $\rho_\eta(z)=\eta^{-1}\rho(z/\eta)$ and
$f_\eta=f*\rho_\eta$. Global Lipschitz continuity gives
\begin{equation}
 \|f_\eta-f\|_\infty\leq L_fm_1\eta,
 \qquad
 \|f_\eta^{(4)}\|_\infty
 \leq L_f\|\rho_\eta'''\|_1=L_f c_3\eta^{-3}.
 \label{eq:cgbm-mollifier-bounds}
\end{equation}
The second estimate uses the weak first derivative $f'$, which exists
almost everywhere and is bounded by $L_f$. Moving the three remaining
derivatives to the smooth kernel proves the formula and its bound.
Since the radius is smaller than the distance of either kink to the
boundary, $f_\eta$ equals $0$ near $L$ and $H$ near $U$.

Subtract the affine lift from $f_\eta$ and shift the domain to
$[0,U-L]$. The lifted function is zero at both endpoints and affine in
their neighbourhoods. Its even derivatives, of order at least two,
vanish there. Odd reflection at both endpoints therefore produces a
smooth function of period $2(U-L)$: the value and every derivative agree
at each reflection. In particular, the fourth derivative of the extension
has maximum norm at most $L_f c_3\eta^{-3}$. Convolve this extension with
the whole-line Gaussian heat kernel and add the affine lift back. Oddness
keeps the lifted solution zero at both boundaries, so this is precisely
the finite-domain smooth solution $u_\eta$. Gaussian convolution contracts
the maximum norm and commutes with derivatives; hence
\begin{equation}
 \sup_{0\leq s\leq\tau}\|\partial_y^4u_\eta(s,\cdot)\|_\infty
 \leq L_f c_3\eta^{-3},\qquad
 \partial_s^2u_\eta=a^2\partial_y^4u_\eta.
 \label{eq:cgbm-smooth-regularity}
\end{equation}

At time $s_n=nk$, insert $u_\eta(s_n,y_j)$ into the numerical equation
and divide its residual by $k$. The integral remainder for backward
time differencing is at most
$k\sup|\partial_s^2u_\eta|/2$. The integral or Taylor remainder for the
centred second difference is at most
$h^2\sup|\partial_y^4u_\eta|/12$. Thus the residual $R_j^n$ obeys
\begin{equation}
 |R_j^n|\leq
 \left(\frac{a^2k}{2}+\frac{ah^2}{12}\right)L_f c_3\eta^{-3}.
 \label{eq:cgbm-residual}
\end{equation}
If the smooth grid starts from $f_\eta$, its error satisfies
$A_h e^n=e^{n-1}-kR^n$ with $e^0=0$. Equation
\eqref{eq:cgbm-mmatrix} and induction give
$\|e^N\|_\infty\leq Nk\sup_n\|R^n\|_\infty$.
Both the continuous maximum principle and the discrete contraction bound
the change from $f_\eta$ back to $f$ by $L_fm_1\eta$. Adding these two
changes proves \eqref{eq:cgbm-be-error}. Taking
$\eta$ proportional to $(h^2+k)^{1/4}$, eventually below the fixed boundary
margin, proves convergence and the stated upper rate. This is a proof
for the actual kinked initial condition, not for a replacement smooth
contract.

\subsection{Domain truncation and interpolation}

Let $Y_s=x+\sigma B_s$, and let $\zeta$ be its first exit from $[L,U]$.
The stopped Brownian representation of the finite-domain solution is
\begin{equation}
 w^{\mathrm{box}}(\tau,x)=
 \mathbb E_x[f(Y_\tau)\mathbf 1_{\{\zeta>\tau\}}]
 +\mathbb E_x[g(Y_\zeta)\mathbf 1_{\{\zeta\leq\tau\}}],
 \qquad g(L)=0,\quad g(U)=H.
 \label{eq:cgbm-stopped-heat}
\end{equation}
It follows either from the stopped It\^o formula for positive-time smooth
solutions and bounded convergence, or from the killed heat semigroup.
The full solution is $\mathbb E_x f(Y_\tau)$. On exit, the terminal payoff
and boundary value are both in $[0,H]$, so their difference is bounded
in absolute value by $H$. The reflection principle bounds each one-sided
exit event by twice its corresponding normal tail; a union bound proves
\eqref{eq:cgbm-domain-beta}. The upper and lower exit contributions can
have opposite signs, so this is an absolute error bound rather than an
ordering of the two prices.

At positive maturity, differentiate the whole-line convolution once using
the bounded weak derivative $f'$, then once using the Gaussian density.
For a Gaussian with standard deviation $v$, the $L^1$ norm of its density
derivative is $\sqrt{2/\pi}/v$. Therefore
\begin{equation}
 \|w_{yy}(\tau,\cdot)\|_\infty
 \leq \frac{L_f}{\sigma\sqrt{\tau}}\sqrt{\frac{2}{\pi}}.
 \label{eq:cgbm-heat-curvature}
\end{equation}
Linear interpolation on an interval of width $h$ has error at most
$h^2\sup|w_{yy}|/8$. Interpolation weights are non-negative and sum to
one. Apply the domain bound at the two bracketing nodes and the nodal
discretisation bound to these same weights; add the whole-line
interpolation error and discount. This proves
\eqref{eq:cgbm-total-error}, without requiring a new boundary regularity
bound for $w^{\mathrm{box}}$.

\subsection{Uniform consistency for Greeks}

For a central first difference, price error contributes at most
$(E+E)/(2d)=E/d$. Taylor's theorem with an integral remainder gives
bias at most $M_3d^2/6$. For the central second difference, price error
contributes at most $(E+2E+E)/d^2=4E/d^2$, and the fourth-order remainder
is at most $M_4d^2/12$. These prove
\eqref{eq:cgbm-first-bump} and \eqref{eq:cgbm-second-bump}.

On a compact parameter set with positive lower bounds for $S$, $\sigma$
and $\tau$, the analytic formula is smooth and its third and fourth
parameter derivatives are bounded. Its transformed queries also lie in
a compact interval. Let $L\to-\infty$ and $U\to\infty$, keep
$h\to0$ and $k\to0$, and use a radius proportional to
$(h^2+k)^{1/4}$. The Gaussian exit tails vanish uniformly because
$\sigma\sqrt{\tau}$ has a finite upper bound, and the other terms in
\eqref{eq:cgbm-total-error} vanish uniformly. The discount factors are
bounded even when $r$ is negative. Thus $E_n\to0$ uniformly. Choosing
the bump sequence as stated in Proposition~\ref{prop:cgbm-bumps} proves
Greek convergence. A fixed domain or a fixed price error does not meet
this argument merely because its bumps are made smaller.

\section{Monte Carlo derivations and uncertainty proofs}
\label{app:cgbm-monte-carlo}

\subsection{Pathwise first derivatives}

The capped payoff is Lipschitz in its asset argument, with derivative
$\mathbf 1_{\{K<x<J\}}$ except at two points. The continuously distributed
$X$ hits those points with probability zero. Its parameter derivatives are
\begin{equation}
 X_S=X/S,\qquad X_\sigma=X(\sqrt{\tau}Z-\sigma\tau),\qquad
 X_r=\tau X,\qquad
 X_\tau=X\left(b+\frac{\sigma Z}{2\sqrt{\tau}}\right).
 \label{eq:cgbm-path-derivatives}
\end{equation}
On any compact parameter neighbourhood with $S,\sigma,\tau>0$, difference
quotients of the Lipschitz payoff are dominated by a constant times
$(1+|Z|)\exp(C|Z|)$, which is integrable against a Gaussian. This follows
from the mean-value theorem for the smooth terminal map $X$ and the
Lipschitz bound for $F$. Dominated convergence therefore justifies
interchanging differentiation and expectation. Differentiating the
discount factor as well gives
\eqref{eq:cgbm-pw-estimators} and \eqref{eq:cgbm-theta-estimator}.
Neither payoff kink contributes a first-order term because the payoff
itself is continuous there.

On the active region, $X\leq J$ and $Z\in(z_-,z_+)$, giving the delta,
vega and theta ranges in \eqref{eq:cgbm-mc-ranges}. For rho the expression
$IX-F(X)$ is exactly $0$, $K$ or $-H$ in the below-strike, active and
capped regions, respectively. This proves its stated tighter range.

\subsection{The likelihood-ratio gamma}

For $x>0$, the lognormal density is
\begin{equation}
 p_S(x)=\frac{1}{xv\sqrt{2\pi}}
 \exp\!\left[-\frac{1}{2}
 \left(\frac{\log(x/S)-b\tau}{v}\right)^2\right].
 \label{eq:cgbm-density}
\end{equation}
Writing $z=(\log(x/S)-b\tau)/v$, direct differentiation gives
\begin{equation}
 \frac{\partial_Sp_S(x)}{p_S(x)}=\frac{z}{Sv},\qquad
 \frac{\partial_S^2p_S(x)}{p_S(x)}
 =\frac{1}{S^2}\left(\frac{z^2-1}{v^2}-\frac{z}{v}\right).
 \label{eq:cgbm-lr-weight}
\end{equation}
On a compact spot neighbourhood away from zero, the density and its first
two spot derivatives admit a common envelope
\begin{equation}
 C(1+|\log x|^2)x^{-1}
 \exp\{-|\log x|^2/(4v^2)\},\qquad x>0,
 \label{eq:cgbm-lr-envelope}
\end{equation}
for a finite constant $C$. Indeed $\log S+b\tau$ is bounded on the
neighbourhood and $(u-m)^2\geq u^2/2-m^2$; the derivative weights are
polynomials of degree at most two in $u=\log x$, with bounded inverse
spot factors. The envelope is integrable by the substitution $u=\log x$,
and remains integrable after multiplication by $F\leq H$.
Two applications of dominated convergence
to $D\int F(x)p_S(x)\,dx$ prove unbiasedness of
\eqref{eq:cgbm-lr-estimator}. Using $F\leq H$ gives
\begin{align}
 \mathbb E[Y_\Gamma^2]
 &\leq \frac{D^2H^2}{S^4}
 \mathbb E\left[
 \left(\frac{Z^2-1}{v^2}-\frac{Z}{v}\right)^2\right]\notag\\
 &=\frac{D^2H^2}{S^4}
 \left(\frac{2}{v^4}+\frac{1}{v^2}\right),
 \label{eq:cgbm-lr-moment}
\end{align}
because $\mathbb E(Z^2-1)^2=2$, $\mathbb E Z^2=1$ and the odd mixed
moment is zero. Variance is no larger than the second moment, proving
\eqref{eq:cgbm-gamma-variance}. Gamma's likelihood-ratio estimator is
unbounded; a bounded-payoff concentration inequality cannot be applied
to it without this separate argument.

\subsection{A bounded mixed gamma rival}

The pathwise delta has the fixed-asset representation
\begin{equation}
 \Delta=\frac{D}{S}\int_K^J x p_S(x)\,dx.
 \label{eq:cgbm-delta-density}
\end{equation}
The integration limits are fixed when spot varies. On this bounded
interval away from zero, $x\partial_S p_S(x)$ has a common integrable
envelope on every compact spot neighbourhood. Differentiation is
therefore legitimate even though the payoff derivative is an indicator.
Using \eqref{eq:cgbm-lr-weight},
\begin{align}
 \partial_S\Delta
 &=-\frac{D}{S^2}\int_K^J x p_S(x)\,dx
 +\frac{D}{S}\int_K^J x\partial_Sp_S(x)\,dx\notag\\
 &=\frac{D}{S^2}\int_K^J x p_S(x)
 \left(\frac{z}{v}-1\right)\,dx.
 \label{eq:cgbm-mixed-gamma-derivation}
\end{align}
This proves unbiasedness of \eqref{eq:cgbm-mixed-gamma}. It includes the
boundary effect through the density derivative. On the active region,
$X<J$ and $|Z|\leq m$, proving \eqref{eq:cgbm-mixed-gamma-bound} and its
finite second moment. As an independent analytic check, Gaussian tilting
reduces its expectation to
$e^{-q\tau}\{\varphi(d_{1,K})-\varphi(d_{1,J})\}/(Sv)$,
in agreement with \eqref{eq:cgbm-delta-gamma}. I combine the established
pathwise and likelihood-ratio constructions of
\citet{broadieglasserman1996}, deriving the capped-payoff expression
explicitly here; I make no priority claim. The ordinary and mixed gamma
estimators are compared using identical financial parameters and matched
terminal evaluation budgets; neither estimator is discarded after
seeing its variance.

\subsection{Confidence guarantees and stronger classical sampling}

For completeness, the bounded-variable concentration argument of
\citet{hoeffding1963} can be checked directly. If $Y$ lies in an interval
of width $R>0$, define
$A(s)=\log\mathbb E\exp\{s(Y-\mathbb E Y)\}$ for real $s$.
Boundedness justifies differentiating twice. The second derivative is
the variance of $Y$ under the exponentially tilted probability law.
That law has the same interval support; centring at its midpoint gives
$A''(s)\leq R^2/4$. Since $A(0)=A'(0)=0$, integration twice yields
\begin{equation}
 A(s)\leq R^2s^2/8,\qquad s\in\mathbb R.
 \label{eq:cgbm-hoeffding-mgf}
\end{equation}
For independent observations and $s>0$, exponential Markov's inequality
and the product moment-generating function therefore give
\begin{equation}
 \mathbb P(\overline Y-\mathbb E Y\geq t)
 \leq\exp\{n(R^2s^2/8-st)\}.
 \label{eq:cgbm-hoeffding-chernoff}
\end{equation}
The minimiser is $s=4t/R^2$, giving the upper-tail bound
$\exp(-2nt^2/R^2)$. Apply the same argument to $-Y$ and sum the two
tail bounds. Solving
$2\exp(-2nt^2/R^2)=\alpha$ proves
\eqref{eq:cgbm-hoeffding}. If $R=0$, the observation is constant and
its confidence radius is zero. Independence gives
$\operatorname{Var}(\overline Y_\Gamma)\leq B_\Gamma^2/n$;
Chebyshev's inequality gives the stated gamma radius. A union bound
requires no independence between different outputs calculated on the
same paths. Their allocated failure probabilities must sum to the
intended familywise level.

For the stock control $C=D X-S e^{-q\tau}$, the Gaussian exponential
moment gives $\mathbb E C=0$ and a finite second moment. For any estimator
$Y$, a deterministic coefficient $\beta$ yields
\begin{equation}
 \operatorname{Var}(Y-\beta C)
 =\operatorname{Var}(Y)-2\beta\operatorname{Cov}(Y,C)
 +\beta^2\operatorname{Var}(C).
 \label{eq:cgbm-control-variance}
\end{equation}
When $\operatorname{Var}(C)>0$, the population minimiser is
$\beta_* =\operatorname{Cov}(Y,C)/\operatorname{Var}(C)$ and the
minimum variance is
$\operatorname{Var}(Y)-\operatorname{Cov}(Y,C)^2/
\operatorname{Var}(C)$. An estimated pilot coefficient need not attain
that minimum. Independence of pilot and main sample preserves conditional
unbiasedness; a zero pilot denominator is handled by selecting zero
control coefficient. Antithetic pairs replace $Y(Z)$ and $C(Z)$ with
their averages at $Z$ and $-Z$ and preserve their expectations, but variance
reduction for every Greek is not guaranteed merely by symmetry. The
comparison therefore reports measured variance and complete sampling work.
For fixed finite pilot coefficients, the main-pair estimators have finite
variance, so their studentised normal intervals have asymptotic validity
under non-zero variance. Finite-sample exact coverage is not claimed.

\section{Conditional projection and control-variate proofs}
\label{app:cgbm-projection}

\subsection{Completeness, closedness and the nearest-point construction}

For this appendix write $\mathbb E=\mathbb E_Q$,
$\langle X,Z\rangle=\mathbb E[XZ]$ and
$\|X\|_2=(\mathbb E[X^2])^{1/2}$. Random variables are identified when
they agree $Q$-almost surely. All variables are real-valued.
First, $L^2(Q)$ is complete. From any Cauchy sequence $(X_n)$ choose a
subsequence $(X_{n_j})$ with
$\|X_{n_{j+1}}-X_{n_j}\|_2\leq2^{-j}$. Minkowski's inequality and
Fatou's lemma give
\begin{equation}
 \left\|\sum_{j\geq k}|X_{n_{j+1}}-X_{n_j}|\right\|_2
 \leq\sum_{j\geq k}2^{-j}.
 \label{eq:cgbm-completeness-tail}
\end{equation}
The non-negative infinite sum is consequently finite almost surely.
The subsequence converges pointwise to a measurable finite limit $X$,
with the value zero assigned on its null non-convergence set. The same
tail bound proves $X\in L^2(Q)$ and $X_{n_j}\to X$ in $L^2$.
The Cauchy property then gives $X_n\to X$ in $L^2$.
If each $X_n$ admits a $\mathcal G$-measurable representative, the
convergence set and pointwise limit of this subsequence are
$\mathcal G$-measurable. Hence $X$ admits a $\mathcal G$-measurable
representative. This proves that $L^2(\mathcal G)$ is closed; linearity
of that space follows directly from measurability and the triangle
inequality.

Fix $Y\in L^2(Q)$ and set
$d=\inf_{Z\in L^2(\mathcal G)}\|Y-Z\|_2$. Choose $(Z_n)$ whose
distances converge to $d$. The midpoint lies in the subspace, so the
parallelogram identity yields
\begin{equation}
 \|Z_n-Z_k\|_2^2
 =2\|Y-Z_n\|_2^2+2\|Y-Z_k\|_2^2
 -4\left\|Y-\frac{Z_n+Z_k}{2}\right\|_2^2
 \leq2\|Y-Z_n\|_2^2+2\|Y-Z_k\|_2^2-4d^2.
 \label{eq:cgbm-projection-existence}
\end{equation}
Thus $(Z_n)$ is Cauchy. Completeness and closedness give a limit
$m\in L^2(\mathcal G)$ attaining $d$. For each
$h\in L^2(\mathcal G)$, the quadratic
\begin{equation}
 \|Y-m-sh\|_2^2
 =\|Y-m\|_2^2-2s\langle Y-m,h\rangle+s^2\|h\|_2^2
 \label{eq:cgbm-projection-orthogonality}
\end{equation}
has its minimum at $s=0$, implying
$\langle Y-m,h\rangle=0$. Conversely this orthogonality proves
\eqref{eq:cgbm-projection-pythagoras}, and therefore minimality.
If two minimisers exist, their difference is in the subspace and
orthogonal to itself, so they agree almost surely.

\subsection{The integral characterisation and conditional Jensen inequality}

Taking $h=\mathbf1_A$ in the orthogonality relation proves the displayed
indicator-integral characterisation in
Theorem~\ref{thm:cgbm-projection}. Conversely, that characterisation
first implies orthogonality against every $\mathcal G$-measurable
simple function by finite linearity. For any $h\in L^2(\mathcal G)$,
truncate $h$ to $[-n,n]$ and then approximate that truncation by a simple
function with pointwise error at most $1/n$. The resulting simple
sequence converges to $h$ in $L^2$, since
$\mathbb E[h^2\mathbf1_{\{|h|>n\}}]\to0$ by dominated convergence.
Cauchy--Schwarz then extends orthogonality to $h$. Thus the integral
characterisation is equivalent to the projection construction.
It is also the usual definition of conditional expectation. Its
uniqueness among integrable $\mathcal G$-measurable candidates follows
by testing the difference on its positive and negative sets.

The conditional expectation of $Y^2$ requires an $L^1$ extension rather
than an assumption that $Y\in L^4$. The projection preserves positivity:
if $U\geq0$, $U\in L^2$, the integral identity on
$A=\{P_{\mathcal G}U<0\}$ contradicts a negative integral unless
$Q(A)=0$. Linearity and this positivity show that the projection
preserves order. For $U\geq0$, $U\in L^1$, define
$h_n=P_{\mathcal G}(U\wedge n)$, choosing simultaneous representatives
with $0\leq h_n\leq h_{n+1}$. Monotone convergence gives
\begin{align}
 h&=\lim_nh_n<\infty\ \text{almost surely},\qquad
 \mathbb E h=\mathbb E U,\notag\\
 \mathbb E[\mathbf1_A h]&=\mathbb E[\mathbf1_A U]\quad(A\in\mathcal G).
 \label{eq:cgbm-conditional-l1}
\end{align}
The almost-sure finiteness follows because $\mathbb E h<\infty$.
Assigning zero on the $\mathcal G$-measurable null set where the limit
is infinite gives an integrable representative. Positive and negative
parts define conditional expectation for every $L^1$ variable. The
integral characterisation proves linearity and positivity of this
extension by uniqueness.

For every $A\in\mathcal G$, $\mathbf1_A m\in L^2(\mathcal G)$, so
orthogonality implies
$\mathbb E[\mathbf1_A Ym]=\mathbb E[\mathbf1_A m^2]$.
Expanding the square on $A$ therefore proves
\begin{equation}
 \mathbb E[(Y-m)^2\mid\mathcal G]
 =\mathbb E[Y^2\mid\mathcal G]-m^2\geq0
 \quad\text{almost surely}.
 \label{eq:cgbm-conditional-square}
\end{equation}
Here products are integrable by Cauchy--Schwarz, and the equality follows
from the $L^1$ integral characterisation on every $A$.
This proves conditional Jensen in the required square case. Taking
expectations proves $\|P_{\mathcal G}Y\|_2\leq\|Y\|_2$.
Linearity follows from the indicator-integral identity and uniqueness;
idempotence follows because the projection fixes every element of its
subspace. Taking $Z=\mathbb E Y$ in
\eqref{eq:cgbm-projection-pythagoras}, and using
$\mathbb E m=\mathbb E Y$, proves the total-variance identity.

\subsection{Nested information and the direction of variance comparisons}

Let $m_i=P_{\mathcal G_i}Y$ for
$\mathcal G_1\subseteq\mathcal G_2$. Every indicator of
$\mathcal G_1$ is also an indicator of $\mathcal G_2$; applying the
integral characterisation twice gives
$P_{\mathcal G_1}m_2=m_1$. Also
$P_{\mathcal G_2}m_1=m_1$ by measurability. The first identity is the
tower property. Orthogonality against $m_2-m_1$ gives
\begin{align}
 \operatorname{Var}(m_2)
 &=\operatorname{Var}(m_1)+\mathbb E[(m_2-m_1)^2],
 \label{eq:cgbm-nested-variance}\\
 \mathbb E[(Y-m_1)^2]
 &=\mathbb E[(Y-m_2)^2]+\mathbb E[(m_2-m_1)^2].
 \label{eq:cgbm-nested-loss}
\end{align}
For the first equality, the cross term vanishes by projecting
$m_2-m_1$ onto $\mathcal G_1$. For the second, $Y-m_2$ is orthogonal
to the $\mathcal G_2$-measurable difference $m_2-m_1$.
These exact identities distinguish dispersion of the conditional mean
from the prediction loss of approximating the terminal variable by it.

\subsection{Multivariate controls, singular covariance and optimal coefficients}

Let $C$ have $d$ centred square-integrable components and let
$\mu=\mathbb E Y$. The matrix $\Sigma_C=\mathbb E[CC^{\mathsf T}]$ is
symmetric and positive semidefinite, since
$v^{\mathsf T}\Sigma_C v=\mathbb E[(v^{\mathsf T}C)^2]\geq0$.
If $v\in\ker \Sigma_C$, this square has zero expectation and
$v^{\mathsf T}C=0$ almost surely. Therefore
$v^{\mathsf T}\kappa_C=\mathbb E[(v^{\mathsf T}C)(Y-\mu)]=0$.
For a finite symmetric matrix, its spectral decomposition gives
$\operatorname{range}\Sigma_C=(\ker \Sigma_C)^\perp$, so $\kappa_C\in\operatorname{range}\Sigma_C$.
Writing $\Sigma_C=U\operatorname{diag}(\lambda_i)U^{\mathsf T}$, define
$\Sigma_C^\dagger=U\operatorname{diag}(d_i)U^{\mathsf T}$, where
$d_i=1/\lambda_i$ when $\lambda_i>0$ and $d_i=0$ when $\lambda_i=0$.
Then $\beta_*=\Sigma_C^\dagger \kappa_C$ satisfies $\Sigma_C\beta_*=\kappa_C$.
Direct square expansion gives
\begin{align}
 \mathcal J(\beta)&=\operatorname{Var}(Y)-2\beta^{\mathsf T}\kappa_C+
 \beta^{\mathsf T}\Sigma_C\beta\notag\\
 &=\operatorname{Var}(Y)-\kappa_C^{\mathsf T}\Sigma_C^\dagger \kappa_C+
 (\beta-\beta_*)^{\mathsf T}\Sigma_C(\beta-\beta_*).
 \label{eq:cgbm-singular-control-minimum}
\end{align}
All minimisers are exactly $\beta_*+\ker \Sigma_C$. A singular coefficient
vector is not uniquely identifiable, but its controlled random variable
is unique almost surely because null directions produce zero controls.

Componentwise conditional projection preserves centring and
square-integrability of $C$. By linearity,
$P_{\mathcal G}(Y-\beta^{\mathsf T}C)=m-\beta^{\mathsf T}C_{\mathcal G}$
for every fixed $\beta$. Its residual
$Y-m-\beta^{\mathsf T}(C-C_{\mathcal G})$ is orthogonal to
$L^2(\mathcal G)$. Applying total variance proves
\eqref{eq:cgbm-joint-projection}. Minimising first on the conditional
side, then evaluating at a raw population minimiser, proves
\begin{equation}
 \min_\beta \mathcal J_{\mathcal G}(\beta)
 \leq \mathcal J_{\mathcal G}(\beta_{\rm raw})
 \leq \mathcal J(\beta_{\rm raw})=\min_\beta \mathcal J(\beta).
 \label{eq:cgbm-oracle-control-order}
\end{equation}
Applying \eqref{eq:cgbm-singular-control-minimum} to the projected
variables proves \eqref{eq:cgbm-coefficient-penalty}, including singular
$\Sigma_{C,\mathcal G}$.

Optimisation does not generally commute with conditioning. For an
explicit example, let $A,B$ be independent symmetric random signs,
$Y=A+B$, $C=A+2B$ and $\mathcal G=\sigma(A)$. Then
$m=A$, $C_{\mathcal G}=A$, $\beta_{\rm raw}=3/5$ and
$\beta_{\mathcal G}=1$. The respective variances are
$\mathcal J(3/5)=1/5$, $\mathcal J_{\mathcal G}(3/5)=4/25$ and
$\mathcal J_{\mathcal G}(1)=0$. The square penalty is exactly $4/25$.
For nested information and projected controls, applying
\eqref{eq:cgbm-nested-variance} to the fixed terminal residual
$Y-\beta^{\mathsf T}C$ shows
$\mathcal J_{\mathcal G_1}(\beta)\leq \mathcal J_{\mathcal G_2}(\beta)$ for every
$\beta$. Taking infima gives the same non-decreasing order for the
population-optimised residual variance.

Conditioning the response while leaving the control untransformed is
a different estimator. Since $m$ and $C_{\mathcal G}$ are
$\mathcal G$-measurable, orthogonality gives
\begin{equation}
 \operatorname{Var}(m-\beta^{\mathsf T}C)
 =\mathcal J_{\mathcal G}(\beta)
 +\beta^{\mathsf T}(\Sigma_C-\Sigma_{C,\mathcal G})\beta.
 \label{eq:cgbm-unprojected-control-penalty}
\end{equation}
The extra term is non-negative. Both estimators are unbiased for a fixed
coefficient, because both controls have mean zero; the distinction here
is a variance penalty rather than a bias assertion.

\subsection{Independent pilots and the limits of the oracle comparison}
\label{app:cgbm-pilot}

Let $\mathcal P$ be the pilot $\sigma$-algebra, independent of the
production pairs $(Y_i,C_i)_{i=1}^n$, which are independent copies of
$(Y,C)$. Choose a $\mathcal P$-measurable coefficient
$\widehat\beta$ with $\|\widehat\beta\|\leq B<\infty$ almost surely.
For example, project a fitted coefficient onto a deterministic Euclidean
ball. This restriction supplies moment control even when an empirical
covariance matrix is nearly singular. Define
\begin{equation}
 \widehat\mu=n^{-1}\sum_{i=1}^n
 (Y_i-\widehat\beta^{\mathsf T}C_i).
 \label{eq:cgbm-independent-pilot-estimator}
\end{equation}
Each summand is square-integrable, because
$(Y_i-\widehat\beta^{\mathsf T}C_i)^2
\leq2Y_i^2+2B^2\|C_i\|^2$.
Independence fixes the production distribution when conditioning on
$\mathcal P$, while centring gives
$\mathbb E[C_i\mid\mathcal P]=0$. Conditional independence therefore
proves
\begin{align}
 \mathbb E[\widehat\mu\mid\mathcal P]&=\mu,
 \label{eq:cgbm-pilot-unbiased}\\
 \operatorname{Var}(\widehat\mu\mid\mathcal P)
 &=n^{-1}\mathcal J(\widehat\beta),\qquad
 \operatorname{Var}(\widehat\mu)
 =n^{-1}\mathbb E[\mathcal J(\widehat\beta)].
 \label{eq:cgbm-pilot-variance}
\end{align}
The unconditional variance follows from total variance; the variance of
the conditional mean is zero. The same argument applies to
$(m_i,C_{\mathcal G,i})$ if the production projections are exact.
For any such pilot,
\begin{equation}
 \operatorname{Var}(\widehat\mu)
 =\frac{\operatorname{Var}(Y)-\kappa_C^{\mathsf T}\Sigma_C^\dagger \kappa_C}{n}
 +\frac{\mathbb E[(\widehat\beta-\beta_*)^{\mathsf T}
 \Sigma_C(\widehat\beta-\beta_*)]}{n}.
 \label{eq:cgbm-pilot-excess-risk}
\end{equation}
This non-negative excess-risk term explains why separately estimated
raw and conditional coefficients need not preserve the oracle ranking.
If the same pilot coefficient is applied to both jointly projected
production variables, however,
\eqref{eq:cgbm-joint-projection} holds conditionally on the pilot and
its population variance order survives averaging over that pilot.
None of these arguments applies automatically to coefficients fitted
on the production sample: mean-zero controls alone do not make a
data-dependent product have expectation zero. Pilot observations remain
part of the computational budget. Normal intervals based on production
sample variance are asymptotic, not finite-sample guarantees. Specifically,
with a fixed finite pilot held constant while $n\to\infty$ and
$\mathcal J(\widehat\beta)>0$ for almost every pilot, the conditional
central limit theorem and consistency of the production sample variance
give a studentised normal limit. The finite second moments proved above
suffice for these two statements; no fourth moment is required.
A zero-variance endpoint is excluded from this studentised limit.

If numerical projections $\widetilde m,\widetilde C_{\mathcal G}$ obey
the almost-sure deterministic bounds
$|\widetilde m-m|\leq\varepsilon_m$ and
$\|\widetilde C_{\mathcal G}-C_{\mathcal G}\|\leq\varepsilon_C$,
then their pilot-controlled mean satisfies
\begin{equation}
 \left|\mathbb E\!\left[
 n^{-1}\sum_i(\widetilde m_i-
 \widehat\beta^{\mathsf T}\widetilde C_{\mathcal G,i})\right]-\mu\right|
 \leq\varepsilon_m+B\varepsilon_C.
 \label{eq:cgbm-projection-evaluator-bias}
\end{equation}
Indeed subtract the exact estimator, apply the triangle and
Cauchy--Schwarz inequalities to each summand, then take expectations.
Without a validated evaluator-error bound, the floating-point computation
does not inherit an exact numerical unbiasedness guarantee from the
projection theorem.

\subsection{GBM conditional valuation, stock controls and endpoint limits}

Under \eqref{eq:cgbm-law}, for $0\leq t<T$,
\begin{equation}
 S_T=S_t\exp\{b(T-t)+\sigma(W_T^Q-W_t^Q)\}.
 \label{eq:cgbm-markov-split}
\end{equation}
The increment is independent of $\mathcal F_t$ and normally distributed
with variance $T-t$. Integrating the bounded payoff against its normal
density gives \eqref{eq:cgbm-conditional-price}. The resulting expression
is measurable with respect to $\sigma(S_t)$. Testing indicators of that
smaller $\sigma$-algebra shows that it is also
$\mathbb E[Y\mid\sigma(S_t)]$. This equality for the terminal payoff
does not imply that the spot-generated $\sigma$-algebras are nested.

The Gaussian exponential moment gives
$\mathbb E[S_T\mid\mathcal F_t]=S_t e^{(r-q)(T-t)}$, proving
\eqref{eq:cgbm-conditional-control} and $\mathbb E c_t=0$.
Moreover
$c_t=S_0e^{-qT}\{\exp(\sigma W_t^Q-\sigma^2t/2)-1\}$, so
\begin{align}
 \operatorname{Var}(C)&=S_0^2e^{-2qT}(e^{\sigma^2T}-1),\notag\\
 \mathbb E[(C-c_t)^2]
 &=S_0^2e^{-2qT}(e^{\sigma^2T}-e^{\sigma^2t}).
 \label{eq:cgbm-control-residual-variance}
\end{align}
The first two moments of the exponential and total variance prove these
expressions and \eqref{eq:cgbm-conditional-control-variance}.
Since $0\leq Y,m_t\leq e^{-rT}H$, continuity of the GBM paths and
continuity of $V$ at positive remaining maturity imply
$m_t\to V(S_0,T)$ in $L^2$ as $t\downarrow0$ by dominated convergence.
Continuity at expiry gives $m_t\to Y$ in $L^2$ as $t\uparrow T$.
Equation~\eqref{eq:cgbm-control-residual-variance} gives the corresponding
$L^2$ limits for the controls. The nested-filtration identities therefore
apply through both endpoint limits. They concern prices under the
fixed model, not parameter derivatives of a simulation coupling; no
conditional Greek identity is asserted here.
\subsection{Analytic raw moments and smooth conditional diagnostics}
\label{app:cgbm-moment-evaluation}

Let $D_T=e^{-rT}$, $J=K+H$, $\theta=\mathbb E Y$ and
$s_C=S_0e^{-qT}$. For $\ell=0,1,2$ and $0\leq A<B\leq\infty$, define
\begin{align}
 M_\ell(A,B)&=\mathbb E[S_T^\ell\mathbf1_{\{A<S_T<B\}}]\notag\\
 &=S_0^\ell e^{\ell bT+\ell^2\sigma^2T/2}
 \left[\Phi\!\left(\frac{\log(B/S_0)-bT-\ell\sigma^2T}{\sigma\sqrt T}\right)
 -\Phi\!\left(\frac{\log(A/S_0)-bT-\ell\sigma^2T}{\sigma\sqrt T}\right)\right].
 \label{eq:cgbm-truncated-moments}
\end{align}
Zero and infinite endpoints use the corresponding normal-tail limits.
To derive the expression, substitute $S_T=S_0e^{bT+\sigma\sqrt T Z}$,
complete the square in $e^{\ell\sigma\sqrt T z}\varphi(z)$ and shift
the normal integration endpoints by $\ell\sigma\sqrt T$. Thus
\begin{align}
 \mathbb E[Y^2]&=D_T^2\{M_2(K,J)-2KM_1(K,J)+K^2M_0(K,J)
 +H^2M_0(J,\infty)\},\label{eq:cgbm-raw-second}\\
 \operatorname{Cov}(Y,C)&=D_T^2\{M_2(K,J)-KM_1(K,J)
 +HM_1(J,\infty)\}-\theta s_C.
 \label{eq:cgbm-raw-stock-covariance}
\end{align}
These follow by splitting the terminal stock into below-strike, active
and capped regions and expanding the respective payoff polynomials.
Subtracting $\theta^2$ gives $\operatorname{Var}(Y)$, while the stock-control
variance is already derived in \eqref{eq:cgbm-control-residual-variance}.

For conditional diagnostics, let $s(z)=S_0e^{bt+\sigma\sqrt t z}$,
$m(z)=e^{-rt}V(s(z),T-t)$ and
$c(z)=e^{-rt-q(T-t)}s(z)-s_C$. Their moment integrals are
\begin{equation}
 v_t=\int_{\mathbb R}(m(z)-\theta)^2\varphi(z)\,dz,
 \qquad \kappa_t=\int_{\mathbb R}(m(z)-\theta)c(z)\varphi(z)\,dz.
 \label{eq:cgbm-conditional-moment-integrals}
\end{equation}
For an independently expressed projection check, condition the residual
$Y-\beta C$ on $S_t=s$. Put $R=T-t$ and
$X_R=s\exp(bR+\sigma\sqrt R Z')$, where $Z'$ is independent standard normal.
The exact inner variance is
\begin{equation}
 \operatorname{Var}(Y-\beta C\mid S_t=s)
 =e^{-2rt}\operatorname{Var}\{e^{-rR}F(X_R)-\beta e^{-rR}X_R\}.
 \label{eq:cgbm-inner-residual-variance}
\end{equation}
The constant $\beta s_C$ does not affect variance. Applying the raw
truncated-moment formulae with spot $s$ and horizon $R$, then integrating
the displayed quantity against the density of $S_t$, evaluates the residual
term in \eqref{eq:cgbm-joint-projection}. This changes the representation
being checked and avoids coarse quadrature of the payoff kinks. Its
floating-point numerical agreement does not certify an integration remainder.

\section{Gaussian estimation and confidence-image proofs}
\label{app:cgbm-estimation}

\subsection{Likelihood, covariance geometry and existence}

I assume the observations in \eqref{eq:cgbm-statistical-model} are
independent and the true $\Sigma$ is strictly positive definite.
For $a\in\mathbb R^d$ and any candidate $A\succ0$, the joint density is
\begin{equation}
 (2\pi\delta)^{-nd/2}(\det A)^{-n/2}
 \exp\left\{-\frac{1}{2\delta}
 \sum_{i=1}^n(R_i-a\delta)^\top A^{-1}(R_i-a\delta)\right\}.
 \label{eq:cgbm-gaussian-likelihood}
\end{equation}
The identity
\begin{equation}
 \sum_i(R_i-a\delta)(R_i-a\delta)^\top
 =C+n(\bar R-a\delta)(\bar R-a\delta)^\top
 \label{eq:cgbm-scatter-decomposition}
\end{equation}
follows by expansion, since $\sum_i(R_i-\bar R)=0$.
For every $A\succ0$, the second term makes the negative log-likelihood
strictly larger unless $a=\bar R/\delta$. Thus this is the unique
maximising mean. After substituting it, minimising the negative
log-likelihood is equivalent to minimising
$g(A)=\log\det A+\operatorname{tr}(A^{-1}B)$, where $B=C/(n\delta)$.
If $B\succ0$, let $D=A^{-1/2}BA^{-1/2}$ with eigenvalues
$\lambda_1,\ldots,\lambda_d>0$. Trace cyclicity and determinant
multiplicativity give
\begin{equation}
 g(A)-g(B)=\operatorname{tr}D-\log\det D-d
 =\sum_{j=1}^d(\lambda_j-\log\lambda_j-1)\geq0.
 \label{eq:cgbm-likelihood-global-optimum}
\end{equation}
The scalar derivative is $1-1/\lambda$, so each summand has its unique
minimum zero at one. Equality therefore requires $D=I$, equivalently
$A=B$. This proves a unique global maximum. The parameter transformation
$\mu=a+q+\operatorname{diag}(A)/2$ is bijective, so the joint
physical-drift/covariance maximiser is unique as well.

If $C$ has rank $k<d$, orthogonal diagonalisation gives
$C=U\operatorname{diag}(c_1,\ldots,c_k,0,\ldots,0)U^\top$, with
$c_j>0$. Choose
$A_\varepsilon=U\operatorname{diag}(c_1/(n\delta),\ldots,
c_k/(n\delta),\varepsilon,\ldots,\varepsilon)U^\top$.
Then $\operatorname{tr}(A_\varepsilon^{-1}B)=k$, while
$\log\det A_\varepsilon$ contains $(d-k)\log\varepsilon$.
As $\varepsilon\downarrow0$, the log-likelihood tends to $+\infty$.
The case $k=0$ uses $A_\varepsilon=\varepsilon I$ and the same argument.
There is no finite maximiser in the open positive-definite parameter
space. The centred $n\times d$ observation matrix has row sum zero,
so $\operatorname{rank}C\leq\min(d,n-1)$. In particular, $n\leq d$
forces this failure. Under the Gaussian model, $C\succ0$ almost surely
when $n>d$: after the rotation below, its factor is a
$(n-1)\times d$ matrix of independent normal entries. A $d\times d$
minor is a nonzero polynomial in continuously distributed entries;
its zero set has Lebesgue measure zero by induction on the number of
variables, using finite univariate root sets and Fubini for nonzero
coefficient polynomials.

\subsection{Orthogonal rotation, exact pivots and drift uncertainty}

Let $Z$ be an $n\times d$ matrix of independent standard normals, and
write the observation matrix as
$R=\delta\mathbf1 a^\top+\sqrt\delta ZL^\top$.
Choose an orthogonal $n\times n$ matrix $O$ whose first row is
$\mathbf1^\top/\sqrt n$. The joint density of $Z$ is proportional to
$\exp(-\|Z\|_F^2/2)$. The transformation $Z\mapsto OZ$ preserves
the Frobenius norm and has absolute Jacobian one. Hence every entry
of $OZ$ is again an independent standard normal. Denote its rows by
$z_0^\top,z_1^\top,\ldots,z_{n-1}^\top$. Taking the first rotated
row and removing that row from the centred matrix yields
\begin{equation}
 \widehat a=a+\frac{Lz_0}{\sqrt{n\delta}},\qquad
 \frac C\delta=L\left(\sum_{i=1}^{\nu}z_i z_i^\top\right)L^\top,
 \qquad \nu=n-1.
 \label{eq:cgbm-normal-wishart-rotation}
\end{equation}
These two quantities are independent. I define
$W_d(\Sigma,\nu)$ as the distribution of
$\sum_{i=1}^{\nu}X_iX_i^\top$ for independent
$X_i\sim N_d(0,\Sigma)$. Equation~\eqref{eq:cgbm-normal-wishart-rotation}
therefore proves $C/\delta\sim W_d(\Sigma,\nu)$, including the singular
case $\nu<d$. Since $\mathbb E[X_iX_i^\top]=\Sigma$,
\begin{equation}
 \mathbb E\widehat\Sigma_{\mathrm U}=\Sigma,
 \qquad \mathbb E\widehat\Sigma_{\mathrm{ML}}=
 \frac{n-1}{n}\Sigma,
 \qquad \operatorname{Cov}(\widehat a)=\frac{\Sigma}{n\delta}.
 \label{eq:cgbm-estimator-moments}
\end{equation}

For one coordinate, put $s_j^2=\widehat\Sigma_{\mathrm U,jj}$.
The $\nu$ projections $(Lz_i)_j/\sqrt{\Sigma_{jj}}$ are independent
standard normals, so
\begin{equation}
 Q_j=\frac{\nu s_j^2}{\Sigma_{jj}}\sim\chi^2_\nu,
 \qquad
 T_j=\frac{\sqrt{n\delta}(\widehat a_j-a_j)}{s_j}
 \sim t_\nu.
 \label{eq:cgbm-exact-pivots}
\end{equation}
Here $\chi^2_\nu$ is a sum of $\nu$ independent squared standard
normals, and $t_\nu$ is the ratio of an independent standard normal
to $\sqrt{Q_j/\nu}$. Both laws are continuous, so their quantiles
give exact marginal coverage. For $0<\alpha<1$, let $x_{\nu,p}$ and
$t_{\nu,p}$ be their respective $p$ quantiles. Solving the central
pivot inequalities for the unknown parameters gives
\begin{align}
 a_j&\in\left[\widehat a_j-t_{\nu,1-\alpha/2}
 \frac{s_j}{\sqrt{n\delta}},\;
 \widehat a_j+t_{\nu,1-\alpha/2}
 \frac{s_j}{\sqrt{n\delta}}\right],
 \label{eq:cgbm-logdrift-ci}\\
 \sigma_j&\in\left[
 \sqrt{\frac{\nu s_j^2}{x_{\nu,1-\alpha/2}}},\;
 \sqrt{\frac{\nu s_j^2}{x_{\nu,\alpha/2}}}\right].
 \label{eq:cgbm-volatility-ci}
\end{align}
The square-root transformation is increasing on positive values;
it changes neither event nor coverage. These are marginal intervals. Assigning error probability
$\alpha/(2d)$ to each of the $d$ mean and $d$ volatility intervals
and applying the union bound gives simultaneous coverage at least
$1-\alpha$, without assuming independence across assets. This covers
no off-diagonal covariance interval. Joint bounds
$a_j\in[a_j^-,a_j^+]$ and $\Sigma_{jj}\in[v_j^-,v_j^+]$
also imply $\mu_j\in[a_j^-+q_j+v_j^-/2,a_j^++q_j+v_j^+/2]$;
thus the same simultaneous event supplies physical-drift bounds.

Gaussian integration by parts gives
$\mathbb E Z^{2m}=(2m-1)\mathbb E Z^{2m-2}$, hence
$\mathbb E Z^2=1$ and $\mathbb E Z^4=3$. More generally, differentiating
$\mathbb E e^{t^\top X}=e^{t^\top\Sigma t/2}$ four times at zero gives
$\mathbb E[X_iX_jX_kX_l]=\Sigma_{ij}\Sigma_{kl}+
\Sigma_{ik}\Sigma_{jl}+\Sigma_{il}\Sigma_{jk}$.
Differentiation is justified by the integrability of a polynomial
in $\|X\|$ times $e^{c\|X\|}$ for each finite $c$, obtained by
completing the Gaussian square. The independent Wishart summands then give
\begin{equation}
 \operatorname{Cov}(\widehat\Sigma_{\mathrm U,ij},
 \widehat\Sigma_{\mathrm U,kl})=
 \frac{\Sigma_{ik}\Sigma_{jl}+\Sigma_{il}\Sigma_{jk}}{\nu},
 \qquad \operatorname{Var}(s_j^2)=\frac{2\Sigma_{jj}^2}{\nu}.
 \label{eq:cgbm-wishart-covariance}
\end{equation}
The physical-drift estimators are
\begin{align}
 \widehat\mu_{\mathrm U,j}&=\widehat a_j+q_j+\tfrac12s_j^2,\\
 \widehat\mu_{\mathrm{ML},j}&=\widehat a_j+q_j+
 \tfrac12\widehat\Sigma_{\mathrm{ML},jj}.
 \label{eq:cgbm-physical-drift-estimators}
\end{align}
Their moments follow from \eqref{eq:cgbm-estimator-moments}:
\begin{align}
 \mathbb E\widehat\mu_{\mathrm U,j}&=\mu_j,\\
 \mathbb E\widehat\mu_{\mathrm{ML},j}&=\mu_j-\frac{\Sigma_{jj}}{2n},\\
 \operatorname{Var}(\widehat\mu_{\mathrm U,j})&=
 \frac{\Sigma_{jj}}{n\delta}+\frac{\Sigma_{jj}^2}{2\nu}.
 \label{eq:cgbm-physical-drift-moments}
\end{align}
The variance uses the mean/scatter independence, not an assumed
zero correlation. The random variance correction makes a Student-$t$
interval for $a_j$ inapplicable directly to $\mu_j$.
For fixed $T_{\mathrm{obs}}=n\delta$, the first variance term is
$\Sigma_{jj}/T_{\mathrm{obs}}$, whereas the covariance estimator's
relative standard deviation is $\sqrt{2/\nu}$. Thus increasingly
frequent observations within this ideal diffusion model leave a
nonzero drift-estimation floor. With fixed $\delta$ and increasing
$n$, these variances tend to zero. Equation~\eqref{eq:cgbm-wishart-covariance}
also makes every covariance-entry variance vanish, proving
mean-square consistency of the covariance matrix, mean and corrected drift
in fixed dimension.

For any $A\succeq0$, $0<\lambda\leq1$ and $\eta>0$, put
$A_\lambda=(1-\lambda)A+\lambda\eta I$. Quadratic forms give
\begin{equation}
 \lambda_{\min}(A_\lambda)\geq\lambda\eta,
 \qquad \operatorname{cond}_2(A_\lambda)
 \leq1+\frac{(1-\lambda)\|A\|_2}{\lambda\eta}.
 \label{eq:cgbm-positive-shrinkage}
\end{equation}
Indeed the smallest eigenvalue is at least $\lambda\eta$ and the
largest is at most $(1-\lambda)\|A\|_2+\lambda\eta$; their ratio
proves the bound. The implemented trace-target rule takes
\begin{equation}
 \widehat\eta=\operatorname{tr}(\widehat\Sigma_{\mathrm{ML}})/d,
 \qquad \widehat\Sigma_{\mathrm{tr}}=
 (1-\lambda)\widehat\Sigma_{\mathrm{ML}}+\lambda\widehat\eta I.
 \label{eq:cgbm-trace-target-rule}
\end{equation}
On positive sample trace, the same deterministic bound holds with
$A=\widehat\Sigma_{\mathrm{ML}}$ and $\eta=\widehat\eta$.
Since $\|A\|_2\leq\operatorname{tr}A$ for $A\succeq0$, the
trace-target condition number is also at most $1+d(1-\lambda)/\lambda$.
The trace is positive almost surely under the model for $n\geq2$;
all observations being identical has probability zero. A zero trace
therefore remains an explicitly rejected deterministic input.
Linearity and \eqref{eq:cgbm-estimator-moments} give
\begin{equation}
 \mathbb E\widehat\Sigma_{\mathrm{tr}}=
 \frac{n-1}{n}\left[(1-\lambda)\Sigma+
 \lambda\frac{\operatorname{tr}\Sigma}{d}I\right].
 \label{eq:cgbm-trace-target-expectation}
\end{equation}
Subtracting $\Sigma$ gives its bias. Consistency of the raw estimator
and continuity of the trace imply convergence to
$(1-\lambda)\Sigma+\lambda\operatorname{tr}(\Sigma)I/d$.
For fixed $\lambda>0$, this differs from $\Sigma$ unless the latter
is spherical. Shrinkage is an established alternative \citep{ledoitwolf2004}.
This rule is neither the unrestricted MLE nor the
Ledoit--Wolf selected shrinkage rule, and its diagonals inherit
no exact Wishart pivot. It is a positivity remedy, not an inference
improvement established by the experiment.

The pricing law uses a different drift. Since $L$ is invertible, set
$\theta=L^{-1}(\mu-r\mathbf1)$. On a finite horizon $T$, the density
$\exp(-\theta^\top W_T^P-\|\theta\|^2T/2)$ has expectation one and
is strictly positive. Completing the squares in the independent
Gaussian increment densities, as in Appendix~\ref{app:cgbm-financial},
makes $W^Q=W^P+\theta t$ Brownian under the equivalent $Q$.
Substitution gives $dS_j/S_j=(r-q_j)dt+(L\,dW^Q)_j$.
The diffusion covariance remains $LL^\top=\Sigma$ in this stipulated
model; this is not evidence that general physical and pricing
volatility parameters coincide empirically.

For completeness, define $\Gamma(x)=\int_0^\infty u^{x-1}e^{-u}\,du$
for $x>0$. Integration by parts gives $\Gamma(x+1)=x\Gamma(x)$.
Normalisation of the $\nu$-dimensional Gaussian in polar coordinates
gives sphere area $2\pi^{\nu/2}/\Gamma(\nu/2)$; substituting
$q=\|z\|^2$ then gives
\begin{equation}
 f_{\chi^2_\nu}(q)=\frac{q^{\nu/2-1}e^{-q/2}}
 {2^{\nu/2}\Gamma(\nu/2)},\qquad q>0.
 \label{eq:cgbm-chi-square-density}
\end{equation}
For the independent normal ratio, substituting $z=t\sqrt{q/\nu}$
with Jacobian $\sqrt{q/\nu}$ into the joint density and integrating
$q$ using the Gamma integral gives
\begin{align}
 f_{t_\nu}(t)&=\frac{1}{\sqrt{2\pi\nu}\,2^{\nu/2}\Gamma(\nu/2)}
 \int_0^\infty q^{(\nu+1)/2-1}
 e^{-q(1+t^2/\nu)/2}\,dq\notag\\
 &=\frac{\Gamma((\nu+1)/2)}{\sqrt{\nu\pi}\Gamma(\nu/2)}
 (1+t^2/\nu)^{-(\nu+1)/2},\qquad t\in\mathbb R.
 \label{eq:cgbm-student-density}
\end{align}

\subsection{The capped price image and coverage}

Let $J=K+H$ and $h=\log(J/K)>0$. Then
$d_{1,K}-d_{1,J}=h/(\sigma\sqrt\tau)>0$, while
\begin{equation}
 d_{1,K}+d_{1,J}
 =\frac{\log(S^2/(KJ))+2(r-q)\tau+\sigma^2\tau}
 {\sigma\sqrt\tau}
 =\frac{\tau(\sigma^2-c_*)}{\sigma\sqrt\tau}.
 \label{eq:cgbm-vega-sign-factorisation}
\end{equation}
Since $\varphi(x)=(2\pi)^{-1/2}e^{-x^2/2}$,
$\varphi(d_{1,K})-\varphi(d_{1,J})$ has the opposite sign to
$d_{1,K}^2-d_{1,J}^2$. Factoring the latter as the difference
times the sum and using the strictly positive difference shows
\begin{equation}
 \operatorname{sgn} V_\sigma=\operatorname{sgn}(c_*-\sigma^2).
 \label{eq:cgbm-vega-sign}
\end{equation}
The prefactor in \eqref{eq:cgbm-vega} is strictly positive.
Thus if $c_*\leq0$ the price is strictly decreasing. If $c_*>0$,
it is strictly increasing below $\sqrt{c_*}$ and strictly decreasing
above it, with exactly one stationary point, a strict global maximum.
There is no interior minimum in either case.

The analytic price is continuous on every compact positive volatility
interval $I=[\ell,u]$. Its range is an interval by the intermediate
value theorem, and its extrema are attained by compactness. The sign
calculation proves
\begin{equation}
 V(I)=\left[\min\{V(\ell),V(u)\},\;
 \max\{V(s):s\in\{\ell,u\}\cup
 \{\sqrt{c_*}:c_*>0,\ \sqrt{c_*}\in I\}\}\right].
 \label{eq:cgbm-confidence-price-image}
\end{equation}
If the stationary point lies strictly between the endpoints, strict
monotonicity gives $V(\sqrt{c_*})>\max\{V(\ell),V(u)\}$.
An endpoint-only upper bound therefore omits a legitimate price.
For any random volatility interval $I(R)$ with
$P\{\sigma\in I(R)\}=1-\alpha$, the deterministic set inclusion
$\{\sigma\in I(R)\}\subseteq\{V(\sigma)\in V(I(R))\}$
proves at least $1-\alpha$ price coverage. Equality need not hold
because the price map may be non-injective. If the true
$\sigma=\sqrt{c_*}$, its strict global maximum belongs to $V(I)$
precisely when $\sigma\in I$, so the corrected coverage is exactly
$1-\alpha$. In contrast, the endpoint-only interval covers this true
price only if an endpoint equals $\sigma$. For the chi-squared interval
in \eqref{eq:cgbm-volatility-ci}, each endpoint is a strictly increasing
continuous function of the positive, continuously distributed $s_j^2$,
so that event has probability zero. This exact counterexample does
not rely on a small numerical discrepancy. The claim assumes all
other price inputs known, a correctly specified sampling law and
the stated equality of diffusion covariance across measures.

Exact pivot coverage does not certify machine calculations of
quantiles or extrema. If an interval is outward-enclosed and each
required price evaluation has proved absolute error at most
$\varepsilon$, subtracting $\varepsilon$ from the smallest computed
endpoint value and adding $\varepsilon$ to the largest computed
candidate value encloses the image, provided the stationary candidate
is included or its uncertainty is enclosed as well. A floating-point
experiment without those bounds supplies a numerical check, not
an exact finite-sample computational certificate.

\subsection{Composition with conditional numerical price guarantees}

In this subsection, $P$ denotes the joint law of the physical data and
the pricing procedure's conditional randomisation, with the stipulated
physical data law as its marginal. The target $V(s)$ remains the
$Q$-expectation in \eqref{eq:cgbm-price}.

Let $D$ denote the data, let $I(D)=[\ell,u]\subset(0,\infty)$ cover the true volatility
with probability at least $1-\alpha$, and let $E(D)=\{\ell,u\}$.
Let $C(D)$ add $\sqrt{c_*}$ when positive and contained in $I(D)$.
The candidates depend only on data and known contractual inputs.
For each candidate, assume a numerical pricing procedure satisfies,
for almost every data realisation,
\begin{equation}
 P\{|\widehat V(s)-V(s)|>\varepsilon_s\mid D\}\leq\beta_s,
 \qquad \sum_{s\in C(D)}\beta_s\leq\beta,
 \quad \varepsilon_s,\beta_s\geq0.
 \label{eq:cgbm-conditional-numerical-confidence}
\end{equation}
The allowances may depend on $D$. A conditional union bound makes
simultaneous numerical success have probability at least $1-\beta$
given $D$, hence also unconditionally by total expectation. Define
\begin{equation}
 \widehat L=\min_{s\in E(D)}(\widehat V(s)-\varepsilon_s),
 \qquad \widehat U=\max_{s\in C(D)}(\widehat V(s)+\varepsilon_s).
 \label{eq:cgbm-composed-price-interval}
\end{equation}
The lower bound uses only endpoints because they contain the true
image minimum; the upper bound uses every extremum candidate.
On simultaneous numerical success,
Equation~\eqref{eq:cgbm-confidence-price-image} gives
$V(I(D))\subseteq[\widehat L,\widehat U]$. Another union bound yields
\begin{equation}
 P\{V(\sigma)\in[\widehat L,\widehat U]\}\geq1-\alpha-\beta.
 \label{eq:cgbm-composed-price-coverage}
\end{equation}
Neither step assumes independence. The certificate requires prices at at most three distinct volatility
values per fixed contract. Classical/quantum comparisons must share interval/candidate information,
declare confidence allowances and charge all pricing calls and estimation
work. The premises
require valid conditional guarantees, including numerical input/evaluation
errors; the floating-point study does not establish them. This standard
composition gives no Greek/basket extension or new quantum algorithm.

\subsection{A validated bridge from rounded observations to price}
\label{app:cgbm-validated-calibration}

I retain the scalar observation law
$X_i\overset{\mathrm{iid}}\sim N(a\delta,\sigma^2\delta)$, $n\geq2$,
with known $\delta>0$, unknown $a$ and $\sigma>0$, and exact declared
$0<\alpha<1$. All other contract inputs are fixed and known. The pricing
target uses their exact supplied rational values; in particular,
$J=K+H$ denotes the exact rational sum and is evaluated outwards. The code observes $z_i$, with proved
rounding bounds $|X_i-z_i|\leq\epsilon_i$. For round-to-nearest binary64,
half the larger adjacent representable gap is a valid radius, including
at zero and for subnormal values; non-finite recorded observations are
rejected. This radius covers one final rounding of an exact return;
returns computed from rounded prices require a separately propagated
observation-error bound. This distinguishes rounded observations from an exactly
Gaussian law on finite binary values.
The target is the exact supplied contract price, rather than a prediction
interval for a future realised payoff. Physical drift is neither needed
for this price certificate nor substituted into the risk-neutral
expectation.

Let $P=I_n-\mathbf1\mathbf1^\top/n$,
$C_X=\|PX\|_2^2$, $C_z=\|Pz\|_2^2$ and
$E^2=\sum_i\epsilon_i^2$. Since $P$ is an orthogonal projection,
$\|P\|_2=1$. The reverse triangle inequality gives
\begin{equation}
 \begin{gathered}
 |\sqrt{C_X}-\sqrt{C_z}|\leq\|P(X-z)\|_2\leq E,\\
 \{\max(0,\sqrt{C_z}-E)\}^2\leq C_X
 \leq(\sqrt{C_z}+E)^2.
 \end{gathered}
 \label{eq:cgbm-rounded-scatter}
\end{equation}
I accumulate $C_z$ and $E^2$ exactly as rational numbers and enclose the
square roots outwards using integer square-root bounds: if
$m=\lfloor 2^b\sqrt{x}\rfloor$, then
$m/2^b\leq\sqrt{x}\leq(m+1)/2^b$. This gives $\underline C\leq C_X\leq\overline C$.
An unresolved positive lower bound is not replaced by an arbitrary floor.

The previously derived pivot is
$Q=C_X/(\sigma^2\delta)\sim\chi^2_\nu$, $\nu=n-1$.
Let $q_L=F_\nu^{-1}(\alpha/2)$ and
$q_U=F_\nu^{-1}(1-\alpha/2)$. I bracket these quantiles by rational
endpoints using outward evaluations of
$F_\nu(x)=\gamma(\nu/2,x/2)/\Gamma(\nu/2)$.
Strict proved inequalities $\overline F_\nu(b_-)<p$
and $\underline F_\nu(b_+)>p$ enclose $F_\nu^{-1}(p)$.
If the retained brackets satisfy $q_L^->0$, then on
$\{q_L\leq Q\leq q_U\}$,
\begin{equation}
 \frac{\underline C}{\delta q_U^+}\leq\sigma^2\leq
 \frac{\overline C}{\delta q_L^-}.
 \label{eq:cgbm-validated-variance}
\end{equation}
Outward square roots give $I=[\ell,u]$. Numerical equality of a CDF
estimate and its target never proves a bracket. Unresolved comparisons
increase precision or produce an explicit abstention at the resource
limit.

For $\ell>0$, continuity on the compact interval $I$ gives attained
extrema; the vega-sign proof already establishes their complete
volatility image. It also gives the branch-free maximiser
\begin{equation}
 s_{\max}=\sqrt{\min\{u^2,\max(\ell^2,c_*)\}},\qquad
 V(I)=[\min\{V(\ell),V(u)\},V(s_{\max})].
 \label{eq:cgbm-clamped-maximum}
\end{equation}
Indeed, $c_*\leq\ell^2$ makes the price non-increasing on $I$,
$c_*\geq u^2$ makes it non-decreasing, and otherwise
$\sqrt{c_*}$ is its unique interior maximum. Exact endpoint coincidences
are included. Clamping both endpoints of an enclosure for $c_*$ and
taking outward square roots therefore encloses $s_{\max}$ without an
unproved membership decision.

I evaluate the call-spread formula with Arb midpoint-radius arithmetic
\citep{johansson2017}, collecting its final term as
$H D\Phi(d_{2,J})$. Inclusion-preserving arithmetic gives endpoint
price enclosures $[L_\ell,U_\ell]$, $[L_u,U_u]$ and a price enclosure
$[L_m,U_m]$ on the maximiser ball. Consequently
\begin{equation}
 \mathcal J=[\min(L_\ell,L_u),U_m]\supseteq V(I).
 \label{eq:cgbm-validated-price-image}
\end{equation}
Cancellation can require more precision but cannot invalidate a proved
enclosure. For fixed $\ell>0$, clamp is Lipschitz, square root is
Lipschitz on $[\ell^2,u^2]$, and the pricing formula is continuous.
Thus refinement shrinks numerical endpoint radii, rather than eliminating
statistical uncertainty; a finite precision budget need not attain the
requested width.

Let $A$ denote successful return of a validated interval. On the pivot
event, whose probability is $1-\alpha$, the rounding and numerical
inclusions imply $V(\sigma)\in\mathcal J$ whenever $A$ occurs. Hence
\begin{equation}
 P\{A,\ V(\sigma)\notin\mathcal J\}\leq\alpha.
 \label{eq:cgbm-validated-return-risk}
\end{equation}
This does not assert $1-\alpha$ coverage conditional on $A$.
Returning the full financial range on every abstention would instead
give unconditional coverage at least $1-\alpha$. Numerical failure
allowance is zero for completed validated enclosures; the sampling and
constant-diffusion hypotheses remain essential. The result does not
certify market-model adequacy, Greeks or quantum advantage.

\begin{thebibliography}{99}
\bibitem[Asmussen(2018)]{asmussen2018}
Asmussen, S. (2018). Conditional Monte Carlo for sums, with applications to insurance and finance. \emph{Annals of Actuarial Science}, 12(2), 455--478. \url{https://doi.org/10.1017/S1748499517000252}. Final publisher PDF consulted, Section 1, printed pages 455--456.

\bibitem[Black and Scholes(1973)]{blackscholes1973}
Black, F. and Scholes, M. (1973). The pricing of options and corporate liabilities. \emph{Journal of Political Economy}, 81(3), 637--654. \url{https://doi.org/10.1086/260062}. Final journal scan consulted, printed pages 637--644.

\bibitem[Broadie and Glasserman(1996)]{broadieglasserman1996}
Broadie, M. and Glasserman, P. (1996). Estimating security price derivatives using simulation. \emph{Management Science}, 42(2), 269--285. \url{https://doi.org/10.1287/mnsc.42.2.269}. Final journal scan consulted, Sections 2--3, printed pages 270--274, and Appendices A--C, printed pages 280--284.

\bibitem[Clopper and Pearson(1934)]{clopperpearson1934}
Clopper, C. J. and Pearson, E. S. (1934). The use of confidence or fiducial limits illustrated in the case of the binomial. \emph{Biometrika}, 26(4), 404--413. \url{https://doi.org/10.1093/biomet/26.4.404}. Final journal scan consulted, Section 1, printed pages 404--407.

\bibitem[Giles and Carter(2006)]{gilescarter2006}
Giles, M. B. and Carter, R. (2006). Convergence analysis of Crank--Nicolson and Rannacher time-marching. \emph{Journal of Computational Finance}, 9(4), 89--112. \url{https://doi.org/10.21314/JCF.2006.152}. Final journal PDF consulted, Sections 1--3, printed pages 89--103, and Section 5, page 111.

\bibitem[Hoeffding(1963)]{hoeffding1963} Hoeffding, W. (1963). Probability inequalities for sums of bounded random variables. \emph{Journal of the American Statistical Association}, 58(301), 13--30. \url{https://doi.org/10.1080/01621459.1963.10500830}. Earlier final-journal-scan consultation retained in the project source ledger, printed pages 13--16, particularly Theorem 2, equation (2.6).

\bibitem[Johansson(2017)]{johansson2017}
Johansson, F. (2017). Arb: efficient arbitrary-precision midpoint-radius interval arithmetic. \emph{IEEE Transactions on Computers}, 66(8), 1281--1292. \url{https://doi.org/10.1109/TC.2017.2690633}. The installed python-flint 0.9.0 interfaces are checked here; the source record distinguishes this API review from the previously cited publication, whose body was not newly consulted.

\bibitem[Ledoit and Wolf(2004)]{ledoitwolf2004}
Ledoit, O. and Wolf, M. (2004). A well-conditioned estimator for large-dimensional covariance matrices. \emph{Journal of Multivariate Analysis}, 88(2), 365--411. \url{https://doi.org/10.1016/S0047-259X(03)00096-4}. Author-hosted publisher-typeset article-in-press version consulted, Sections 1 and 2.1, printed pages 365--369.

\bibitem[Merton(1973)]{merton1973}
Merton, R. C. (1973). Theory of rational option pricing. \emph{The Bell Journal of Economics and Management Science}, 4(1), 141--183. \url{https://doi.org/10.2307/3003143}. Final journal PDF consulted, selected material in Sections 6--7, printed pages 162--166 and 168--171, particularly equation (44), page 170, and footnote 62, page 171.

\bibitem[Merton(1980)]{merton1980}
Merton, R. C. (1980). On estimating the expected return on the market. \emph{Journal of Financial Economics}, 8(4), 323--361. \url{https://doi.org/10.1016/0304-405X(80)90007-0}. Published successor metadata verified; the consulted text is NBER Working Paper 444, February 1980, \emph{On estimating the expected return on the market: an exploratory investigation}, Introduction, pages 4--5, and Appendix A, pages A1--A4, \url{https://www.nber.org/system/files/working_papers/w0444/w0444.pdf}. The final journal text was not consulted.

\bibitem[Nelson(1990)]{nelson1990}
Nelson, B. L. (1990). Control variate remedies. \emph{Operations Research}, 38(6), 974--992. \url{https://doi.org/10.1287/opre.38.6.974}. Final author-hosted journal scan consulted, Sections 1--2, printed pages 975--977, and Section 6, pages 979--980.

\bibitem[Student(1908)]{student1908}
Student (1908). The probable error of a mean. \emph{Biometrika}, 6(1), 1--25. \url{https://doi.org/10.1093/biomet/6.1.1}. Final journal scan consulted, Introduction and Sections I--III, printed pages 1--8.

\bibitem[Wishart(1928)]{wishart1928}
Wishart, J. (1928). The generalised product moment distribution in samples from a normal multivariate population. \emph{Biometrika}, 20A(1--2), 32--52. \url{https://doi.org/10.1093/biomet/20a.1-2.32}. Final journal scan consulted, Introduction and Section 3, printed pages 32--34 and 38--40, especially equations (9)--(12).
\end{thebibliography}
\end{document}
